% ============================================================
%  Bayesian Statistics and Photonics:
%  Failures of Traditional Models and Solutions via the
%  RSVP Framework and the Spherepop/KES Calculus
%
%  Flyxion, Independent Researcher
% ============================================================
\documentclass[11pt, openright, twoside]{book}

% ── Fonts & encoding ─────────────────────────────────────────
\usepackage[T1]{fontenc}
\usepackage{microtype}

% ── Mathematics ───────────────────────────────────────────────
\usepackage{amsmath, amssymb, amsthm, mathtools}
\usepackage{bm}

% ── Page geometry ─────────────────────────────────────────────
\usepackage[
  paperwidth=6.5in, paperheight=9.5in,
  top=1in, bottom=1in,
  inner=1.1in, outer=0.9in,
  includeheadfoot
]{geometry}

% ── Headers & footers ─────────────────────────────────────────
\usepackage{fancyhdr}
\pagestyle{fancy}
\fancyhf{}
\fancyhead[LE]{\small\itshape\leftmark}
\fancyhead[RO]{\small\itshape\rightmark}
\fancyfoot[C]{\small\thepage}
\renewcommand{\headrulewidth}{0.4pt}

% ── Chapter style ─────────────────────────────────────────────
\usepackage{titlesec}
\titleformat{\chapter}[display]
  {\normalfont\Large\bfseries}
  {\chaptername\ \thechapter}{16pt}{\Huge}
\titlespacing*{\chapter}{0pt}{50pt}{40pt}

% ── Boxes ─────────────────────────────────────────────────────
\usepackage{tcolorbox}
\tcbuselibrary{breakable, skins}

\newtcolorbox{implbox}[1][]{
  colback=gray!8, colframe=gray!50,
  fonttitle=\bfseries\small, title=Technical Note,
  breakable, #1}

\newtcolorbox{epistbox}[1][]{
  colback=blue!5, colframe=blue!30,
  fonttitle=\bfseries\small, title=Epistemic Status,
  breakable, #1}

\newtcolorbox{exbox}[1][]{
  colback=green!5, colframe=green!40,
  fonttitle=\bfseries\small, title=Worked Example,
  breakable, #1}

\newtcolorbox{compbox}[1][]{
  colback=red!3, colframe=red!40,
  fonttitle=\bfseries\small, title=Orthodox vs.\ RSVP,
  breakable, #1}

\newtcolorbox{depbox}[1][]{
  colback=yellow!4, colframe=orange!40,
  fonttitle=\bfseries\small, title=Chapter Dependencies,
  breakable, #1}

% ── Hyperlinks ────────────────────────────────────────────────
\usepackage[hidelinks]{hyperref}

% ── Graphics & tables ─────────────────────────────────────────
\usepackage{graphicx, booktabs, array, longtable}

% ── Algorithms ────────────────────────────────────────────────
\usepackage{algorithm, algpseudocode}

% ── Index ─────────────────────────────────────────────────────
\usepackage{makeidx}
\makeindex

% ── Bibliography ─────────────────────────────────────────────
\usepackage[authoryear, round]{natbib}

% ── Theorem environments ──────────────────────────────────────
\theoremstyle{plain}
\newtheorem{theorem}{Theorem}[chapter]
\newtheorem{lemma}[theorem]{Lemma}
\newtheorem{proposition}[theorem]{Proposition}
\newtheorem{corollary}[theorem]{Corollary}

\theoremstyle{definition}
\newtheorem{definition}[theorem]{Definition}
\newtheorem{example}[theorem]{Example}
\newtheorem{axiom}[theorem]{Axiom}

\theoremstyle{remark}
\newtheorem{remark}[theorem]{Remark}
\newtheorem{conjecture}[theorem]{Conjecture}

% ── Notation ─────────────────────────────────────────────────
\newcommand{\vfield}{\mathbf{v}}
\newcommand{\Sfield}{S}
\newcommand{\rsvp}{(\Phi,\vfield,\Sfield)}
\newcommand{\Pop}{\mathbf{Pop}}
\newcommand{\Bind}{\mathbf{Bind}}
\newcommand{\Collapse}{\mathbf{Collapse}}
\newcommand{\Refuse}{\mathbf{Refuse}}
\newcommand{\KES}{\Omega_t \to H_{t+1}}
\newcommand{\E}[1]{\mathbb{E}\!\left[#1\right]}
\newcommand{\given}{\,|\,}
\newcommand{\Normal}{\mathcal{N}}
\newcommand{\Admiss}{\mathcal{A}}
\newcommand{\Traj}{X}
\newcommand{\Rep}{M}
\newcommand{\proj}{\pi}
\newcommand{\kvec}{\mathbf{k}}
\newcommand{\xvec}{\mathbf{x}}

\setlength{\parskip}{4pt}
\setlength{\parindent}{1.4em}

% ─────────────────────────────────────────────────────────────
\begin{document}
\frontmatter

% ── Title page ───────────────────────────────────────────────
\begin{titlepage}
\centering
\vspace*{2cm}
{\Huge\bfseries Bayesian Statistics and Photonics\par}
\vspace{0.6cm}
{\Large Failures of Traditional Models and Solutions\\[4pt]
via the RSVP Framework\\[4pt]
and the Spherepop/KES Calculus\par}
\vspace{2cm}
{\large Flyxion\par}
{\normalsize Independent Researcher\par}
\vfill
\end{titlepage}
\clearpage{\pagestyle{empty}\cleardoublepage}

% ── Preface ──────────────────────────────────────────────────
\chapter*{Preface}
\addcontentsline{toc}{chapter}{Preface}

This textbook undertakes a foundational reconstruction of Bayesian inference and quantum
photonics. It begins from the claim that existing frameworks operate on the wrong primitive
objects. Standard Bayesian inference takes probability distributions over predefined
hypothesis spaces as its primary objects. Standard quantum optics takes Hilbert space
vectors and the Born rule as its primary objects. Both are compressions of a richer
underlying structure: the admissible trajectory field from which observable quantities
and statistical distributions both emerge.

The book is organized around four structural ideas that recur at every level of the
exposition. The first is \emph{admissibility}: not all configurations are equally
accessible at a given time, and this differential accessibility is the true primitive from
which probability distributions emerge. The second is \emph{projection}: the map from rich
trajectory space to compressed representational manifold, whose fiber structure encodes the
information lost by conventional models. The third is \emph{entropy as accessibility}: the
entropy field $\Sfield$ measures future trajectory volume, not merely Shannon uncertainty.
The fourth is \emph{irreversible event structure}: the Spherepop/KES calculus treats events
as primitive and derives states as equivalence classes of event histories, placing
irreversibility at the foundation rather than deriving it from reversible dynamics.

Every chapter of this text is anchored to one or more of these four ideas. The reader who
loses the thread can always return to them.

\medskip
The book is organized as follows. Part~0 covers mathematical prerequisites: Hilbert spaces,
manifolds, Fisher-Rao geometry, and rigorous Bayesian inference. Part~I is diagnostic,
identifying specific failures of orthodox models. Part~II develops the RSVP framework as
a field-theoretic replacement, beginning with admissibility geometry as the conceptual
foundation. Part~III develops the Spherepop/KES irreversible event calculus. Part~IV
unifies the two frameworks, compares them to existing programs, and catalogues experimental
predictions. Part~V develops extensions: projection failure, information geometry beyond
Fisher-Rao, non-Markovian dynamics, photonics as constraint propagation, topological
defects, thermodynamics, learning, and cosmological optics. Part~VI closes with a unified
accessibility calculus and open problems.

\medskip
\textbf{Epistemic status labels} used throughout:
\begin{description}
  \item[Theorem/Proposition] Mathematically derived from stated assumptions.
  \item[Conjecture] Believed true but not yet proved within this framework.
  \item[Ansatz] A working assumption adopted for tractability.
  \item[Interpretation] Physical or inferential reading of a formal result.
  \item[Prediction] Empirically testable consequence, dependent on stated ansätze.
\end{description}

\bigskip
\noindent\textit{Flyxion, Independent Researcher}

% ── Notation ─────────────────────────────────────────────────
\chapter*{Notation and Symbols}
\addcontentsline{toc}{chapter}{Notation and Symbols}

\begin{longtable}{@{}lll@{}}
\toprule
\textbf{Symbol} & \textbf{Name} & \textbf{Interpretation} \\
\midrule\endhead
$\Phi$ & Scalar field & Likelihood potential / salience landscape \\
$\vfield$ & Vector field & Admissibility flow / information transport \\
$\Sfield$ & Entropy field & Local accessibility density \\
$\rsvp$ & RSVP triple & The full RSVP field configuration \\
$\mathcal{M}$ & Plenum manifold & Parameter / state space \\
$g$ & Riemannian metric & Geometry of the plenum \\
$\Traj$ & Trajectory space & Space of all admissible histories \\
$\Rep$ & Representational manifold & Compressed model space \\
$\proj:\Traj\to\Rep$ & Projection & Compression map \\
$\proj^{-1}(m)$ & Fiber & Trajectories projecting to $m\in\Rep$ \\
$\Admiss(x,t)$ & Admissibility & Dynamic accessibility of $x$ at time $t$ \\
$\Omega_t$ & World-state & Constraint-weighted distribution \\
$H_t$ & History state & Realized causal event sequence \\
$\Pop$ & Pop operator & Event to immediate consequence \\
$\Bind$ & Bind operator & Joins causally compatible events \\
$\Collapse$ & Collapse operator & Identifies compatible configurations \\
$\Refuse$ & Refuse operator & Filters events by consistency \\
$\mathcal{F}(e)$ & KES functional & Selection weight for event $e$ \\
$\pi_{\rm RSVP}$ & RSVP prior & Dynamically generated prior \\
$\rho$ & Density operator & Quantum optical state \\
$E_k$ & POVM element & Measurement operator \\
$p_k$ & Born probability & $\operatorname{tr}(\rho E_k)$ \\
$S(\rho)$ & von Neumann entropy & $-\operatorname{tr}(\rho\log\rho)$ \\
$F_{ij}$ & Fisher information & $\mathbb{E}[\partial_i\log\mathcal{L}\,\partial_j\log\mathcal{L}]$ \\
$g^{\rm FR}$ & Fisher-Rao metric & Information geometry metric \\
$g^{\rm RSVP}$ & RSVP metric & Extended accessibility metric \\
$D(\sigma_1,\sigma_2)$ & Trace distance & $\frac{1}{2}\|\sigma_1-\sigma_2\|_1$ \\
$\Delta$ & Laplace-Beltrami & Geometric Laplacian on $(\mathcal{M},g)$ \\
$\hat{n}$ & Photon number & $a^\dagger a$ \\
$Q$ & Mandel parameter & Sub/super-Poissonian character \\
$\Gamma(m)$ & Degeneracy ratio & Fiber volume relative to total \\
$\Gamma_D$ & Decoherence rate & Rate of distinguishability loss \\
$K_m$ & KES memory kernel & $m$-step non-Markovian kernel \\
$\mathfrak{A}$ & Total admissibility & $\int_\Traj\Admiss(x,t)\,dx$ \\
\bottomrule
\end{longtable}

\tableofcontents
\mainmatter

% =============================================================
\part{Mathematical Prerequisites}
% =============================================================

\chapter{Hilbert Spaces, Operators, and Quantum Optical States}

\section{Hilbert Spaces and Fock Space}

A \emph{Hilbert space} $\mathcal{H}$ is a complete inner product space over $\mathbb{C}$.
The bosonic Fock space for quantum optics is
$\mathcal{F}(\mathcal{H}_1) = \bigoplus_{n=0}^\infty \mathcal{H}_1^{\otimes_s n}$,
where $\mathcal{H}_1 = L^2(\mathbb{R}^3)\otimes\mathbb{C}^2$ is the single-photon Hilbert
space and $\otimes_s$ denotes the symmetrized tensor product.

\begin{definition}[Density operator]
A \emph{density operator} $\rho:\mathcal{H}\to\mathcal{H}$ satisfies $\rho=\rho^\dagger$,
$\rho\geq 0$, and $\operatorname{tr}(\rho)=1$. Pure states satisfy $\rho^2=\rho$.
\end{definition}

The creation and annihilation operators satisfy $[a_k,a_{k'}^\dagger]=\delta_{kk'}$.
Coherent states $|\alpha\rangle = e^{-|\alpha|^2/2}\sum_n \frac{\alpha^n}{\sqrt{n!}}|n\rangle$
satisfy $a|\alpha\rangle=\alpha|\alpha\rangle$ and are overcomplete. The Wigner function
$W(\alpha) = \frac{2}{\pi}\operatorname{tr}(\rho\,D(-\alpha)(-1)^{a^\dagger a}D(\alpha))$
is a quasi-probability distribution whose negativity witnesses non-classicality.

\section{Quantum Optical Measurements}

\begin{definition}[POVM]
A \emph{positive operator-valued measure} is a set $\{E_k\}$ with $E_k\geq 0$ and
$\sum_k E_k=\mathbf{1}$. The probability of outcome $k$ is $p_k=\operatorname{tr}(\rho E_k)$.
\end{definition}

POVMs generalize projection-valued measures to realistic detectors that never project onto
sharp eigenstates. Informationally complete POVMs with $M\geq d^2$ elements are required
for quantum state tomography of a $d$-dimensional system.

\chapter{Manifolds, Riemannian Geometry, and Statistical Manifolds}

\section{Smooth Manifolds and Differential Operators}

An $n$-dimensional smooth manifold $\mathcal{M}$ is locally homeomorphic to $\mathbb{R}^n$
with smooth transition maps. A Riemannian metric $g$ assigns an inner product to each
tangent space $T_p\mathcal{M}$. The key operators used throughout are: gradient $\nabla f$,
divergence $\operatorname{div}(X) = \frac{1}{\sqrt{\det g}}\partial_\mu(\sqrt{\det g}\,X^\mu)$,
and Laplace-Beltrami $\Delta f = \operatorname{div}(\nabla f)$.

\section{Statistical Manifolds and the Fisher-Rao Metric}

A \emph{statistical manifold} parameterizes a family of probability distributions
$\{p(\cdot\given\theta)\}_{\theta\in\mathcal{M}}$. The Fisher-Rao metric is
\begin{equation}
  g^{\rm FR}_{\mu\nu}(\theta) = \int p(x\given\theta)\,
  \frac{\partial\log p}{\partial\theta^\mu}\frac{\partial\log p}{\partial\theta^\nu}\,dx.
\end{equation}

\begin{theorem}[Chentsov, 1972]
Up to a constant, the Fisher-Rao metric is the unique Riemannian metric on the manifold of
probability distributions invariant under Markov morphisms.
\end{theorem}

The RSVP metric extends the Fisher-Rao metric by incorporating entropy gradients (Chapter~17).

\section{Entropy and Variational Principles}

The Jaynes maximum entropy principle selects the distribution $p^*$ maximizing
$H(p)=-\int p\log p$ subject to moment constraints $\mathbb{E}_p[f_k]=c_k$. The solution
is always an exponential family $p^*(x)=Z^{-1}e^{-\sum_k\lambda_k f_k(x)}$. The RSVP
framework implements this principle \emph{dynamically} via the field equation for
$\Sfield$ rather than as a static optimization.

\chapter{Bayesian Inference: A Rigorous Foundation}

\section{Probability Spaces and Bayes' Theorem}

A probability space $(\Omega,\mathcal{F},P)$ consists of a sample space, $\sigma$-algebra,
and countably additive measure with $P(\Omega)=1$. Bayes' theorem in measure-theoretic form
is
\begin{equation}
  \pi(\theta\given x) = \frac{p(x\given\theta)\,\pi(\theta)}
  {\int p(x\given\theta')\,\pi(\theta')\,d\nu(\theta')},
  \label{eq:bayes}
\end{equation}
valid when the posterior is proper ($\int\pi(\theta\given x)\,d\nu=1$) and $p(x)>0$.

\section{Conjugate Priors and Exponential Families}

A prior $\pi(\theta)$ is \emph{conjugate} to $p(x\given\theta)$ if the posterior lies in
the same parametric family. For the Poisson likelihood relevant to photon counting, the
conjugate prior is $\Gamma(a,b)$ and the posterior given counts $\{n_i\}$ is
$\Gamma(a+\sum n_i, b+N)$, with posterior mean $(a+\sum n_i)/(b+N)$.

\section{Variational Inference and the ELBO}

When the posterior is intractable, variational inference maximizes the evidence lower bound
\begin{equation}
  \mathrm{ELBO}(\phi) = \mathbb{E}_{q_\phi}[\log p(x\given\theta)] - \mathrm{KL}(q_\phi\|\pi).
  \label{eq:elbo}
\end{equation}
The ELBO is tight only when $q_\phi = \pi(\cdot\given x)$. The RSVP framework avoids the
structural assumptions of variational families by treating the prior as a dynamically
evolving field rather than a fixed parametric form.

% =============================================================
\part{Failures of Traditional Models}
% =============================================================

\chapter{The Born Rule and Its Discontents}

\begin{depbox}
This chapter uses Definition 1.1 (density operator), Definition 1.2 (POVM), and
Hilbert space theory from Chapter~1.
\end{depbox}

\section{The Measurement Problem}

The Born rule $p_k = \operatorname{tr}(\rho E_k)$ is a postulate, not a theorem. The
Schr\"{o}dinger equation is linear, deterministic, and unitary; measurement outcomes are
discrete, stochastic, and irreversible. No continuous deformation of the unitary dynamics
produces the Born rule.

\begin{theorem}[Gleason, 1957]
In a separable Hilbert space of dimension $\geq 3$, every frame function
$\mu:\mathcal{P}(\mathcal{H})\to[0,1]$ with $\mu(\mathcal{H})=1$ arises from a unique
density operator $\rho$ via $\mu(V)=\operatorname{tr}(\rho P_V)$.
\end{theorem}

Gleason's theorem shows the Born rule is the \emph{only} consistent assignment; it does
not explain why measurement outcomes should form a frame function. The theorem also fails
in dimension two, precisely where single-qubit polarization states live.

\begin{compbox}
\textbf{Orthodox assumption:} The Born rule $p_k=\operatorname{tr}(\rho E_k)$ is a
fundamental postulate appended to the unitary dynamics.

\textbf{RSVP replacement:} Born statistics emerge as a limiting case of the KES selection
functional (Theorem~\ref{thm:kes_born}) when RSVP fields are trivial ($\Sfield=\mathrm{const}$,
$\Phi=\log p$). In non-trivial backgrounds, the effective probability acquires field-dependent
corrections.

\textbf{Orthodox pathology:} No dynamical explanation for measurement irreversibility.

\textbf{RSVP resolution:} Irreversibility is primitive in the Spherepop event calculus;
the Schr\"{o}dinger equation is a limit, not the foundation.
\end{compbox}

\section{QBism and the Epistemic Deflation}
\label{sec:qbism}

QBism treats $\rho$ as an agent's personal probability assignment, and the Born rule as a
consistency norm analogous to Dutch book coherence. The SIC-POVM reformulation gives
\begin{equation}
  p(E_k) = (d+1)\sum_j q_j r(E_k|\pi_j) - 1,
\end{equation}
departing from the classical law of total probability by an additive correction. This
correction has no dynamical interpretation in QBism. The RSVP framework provides a
field-theoretic account: the correction term is the difference between RSVP-weighted
and classical probability sums arising from non-trivial $\Sfield$ geometry.

\section{The Path Integral and Its Measure Problem}

The Feynman path integral $\mathcal{A} = \int\mathcal{D}[\gamma]\,e^{iS[\gamma]/\hbar}$
is the closest existing approach to admissibility geometry. The RSVP framework replaces
the oscillatory weight $e^{iS/\hbar}$ with the real, positive admissibility weight
$\Admiss(x,t)\propto\Sfield\,e^\Phi$, avoiding the phase-cancellation problem of
the Feynman integral while retaining the sum-over-histories structure.

\begin{exbox}
\textbf{Single-photon polarimetry.} A photon with state $\rho=\mathbf{1}/2$ (maximally
mixed) measured by a polarizing beam splitter gives $p_H=p_V=1/2$ by Born. In the RSVP
framework, this equal probability arises from a flat entropy field $\Sfield=\mathrm{const}$
on the Bloch sphere. When $\Sfield$ is non-uniform due to prior optical interactions, the
effective probabilities shift by an $O(\alpha\Sfield)$ correction testable via photon
statistics.
\end{exbox}

\chapter{Prior Specification in Quantum Optical Inference}

\begin{depbox}
This chapter depends on Definition 3.1 (Jeffreys prior), Bayes' theorem~\eqref{eq:bayes},
and Section~3.2 (conjugate priors).
\end{depbox}

\section{The Prior Assignment Problem}

No canonical rule exists for specifying $\pi(\theta)$ prior to data. The Jeffreys prior
$\pi_J(\theta)\propto\sqrt{\det F(\theta)}$ is reparameterization-invariant but concentrates
mass near the boundary of the density operator manifold for mixed-state tomography, and is
computationally intractable for $n\geq 3$ qubits.

\begin{proposition}[Positivity violation under Gaussian priors]
\label{prop:gaussian_failure}
Let $\tilde\rho$ be drawn from a Gaussian prior on the entries of a $d\times d$ Hermitian
matrix with variance $\sigma^2>0$. Then $\Pr(\tilde\rho\geq 0)\to 0$ as $d\to\infty$.
\end{proposition}

\begin{proof}
The eigenvalues of a GUE matrix with variance $\sigma^2$ follow the Wigner semicircle law
on $[-2\sigma,2\sigma]$, which assigns positive mass to $(-2\sigma,0)$. By Soshnikov's
edge universality (1999), $\lambda_{\min}$ has Tracy-Widom fluctuations around $-2\sigma$,
so $\Pr(\lambda_{\min}\geq 0)\to 0$ as $d\to\infty$.
\end{proof}

\begin{compbox}
\textbf{Orthodox assumption:} Gaussian or Jeffreys priors are standard choices for quantum
state tomography, enforced by Cholesky reparameterization.

\textbf{RSVP replacement:} The RSVP prior $\pi_{\rm RSVP}(\theta)=\Sfield(\theta)e^{\Phi(\theta)}/Z$
is strictly positive by Theorem~\ref{thm:rsvp_prior_well}, derived from field dynamics
rather than algebraic constraint.

\textbf{Orthodox pathology:} Positivity must be enforced post hoc, introducing Jacobian
biases toward pure states.

\textbf{RSVP resolution:} Positivity is a dynamical consequence of $\Sfield>0$.
\end{compbox}

\section{Maximum Entropy Priors and the UV Divergence}

The Jaynes maximum entropy multimode prior
$Z = \prod_k(1-e^{-\beta_k})^{-1}$ diverges unless an ultraviolet cutoff is imposed.
The RSVP entropy field resolves this via Theorem~\ref{thm:uv_finite}: the coupling
structure of the field equations provides intrinsic UV regularization.

\chapter{Coherence, Decoherence, and the Entropy of Observation}

\begin{depbox}
This chapter depends on the Lindblad equation~\eqref{eq:lindblad}, the decoherence
functional (Definition~\ref{def:decoherence_functional}), and Section~2.1 (manifolds).
\end{depbox}

\section{The Lindblad Master Equation and Its Limits}

\begin{equation}
  \frac{d\rho}{dt} = -\frac{i}{\hbar}[H,\rho]
  + \sum_k\gamma_k\!\left(L_k\rho L_k^\dagger - \tfrac{1}{2}\{L_k^\dagger L_k,\rho\}\right).
  \label{eq:lindblad}
\end{equation}

The pathology is not in the Lindblad form but in the post hoc assignment of $L_k$ and
$\gamma_k$: these are fit to data, converting a dynamical equation into a statistical
interpolation device with no principled connection to the Bayesian update on $\rho$.

\begin{definition}[Decoherence functional]
\label{def:decoherence_functional}
For bipartite system-environment evolution $U_t$ from $|\psi_0\rangle_S\otimes|\varepsilon_0\rangle_E$,
the decoherence functional for histories $\alpha,\alpha'$ is
$D(\alpha,\alpha') = \operatorname{tr}(C_\alpha\rho_0 C_{\alpha'}^\dagger)$
where $C_\alpha = P_{\alpha_n}^{t_n}U_{t_n-t_{n-1}}\cdots P_{\alpha_1}^{t_1}U_{t_1}$.
Histories are \emph{consistent} if $D(\alpha,\alpha')\approx 0$ for $\alpha\neq\alpha'$.
\end{definition}

The pointer basis problem: the preferred basis for consistency is not selected by the
formalism. The Spherepop $\Refuse$ operator resolves this (Theorem~\ref{thm:refuse_pointer}).

\begin{theorem}[Entropy non-monotonicity under non-Markovian evolution]
\label{thm:non_markov_entropy}
Let $\{\mathcal{E}_t\}$ be a non-divisible family of CPTP maps. Then there exist initial
states $\rho_1,\rho_2$ such that $D(\mathcal{E}_t(\rho_1),\mathcal{E}_t(\rho_2))$ is
non-monotone in $t$.
\end{theorem}

\begin{proof}
Non-divisibility means $\mathcal{E}_{t,s}=\mathcal{E}_t\circ\mathcal{E}_s^{-1}$ is not
completely positive for some $t>s$. By the Choi-Jamio\l{}kowski isomorphism, there exists
a bipartite state $\tau$ such that $(\mathcal{E}_{t,s}\otimes\mathrm{id})(\tau)$ has a
negative eigenvalue. Constructing $\rho_1,\rho_2$ from its eigenvectors produces a
backflow of distinguishability, violating monotone decay.
\end{proof}

\chapter{Statistical Pathologies in Quantum Optical Tomography}

\begin{depbox}
This chapter depends on the MLE objective~\eqref{eq:mle_tomography}, Proposition~\ref{prop:gaussian_failure},
and the density operator manifold structure from Chapter~1.
\end{depbox}

\section{Boundary Concentration of MLE Estimates}

\begin{equation}
  \hat\rho_{\rm MLE} = \arg\max_{\rho\geq 0,\,\operatorname{tr}\rho=1}
  \sum_k n_k\log\operatorname{tr}(\rho E_k).
  \label{eq:mle_tomography}
\end{equation}

\begin{proposition}[Boundary concentration]
\label{prop:mle_boundary}
Let $\operatorname{rank}(\rho^\star)=r<d$. For finite $N$, $\hat\rho_{\rm MLE}$ is almost
surely rank-$d$, and Fisher-information confidence regions at $\hat\rho_{\rm MLE}$ are
asymptotically incorrect because $\rho^\star$ lies on the boundary of the density matrix manifold.
\end{proposition}

\begin{proof}
The log-likelihood is strictly concave on the interior of the density matrix body, with
gradient diverging to $-\infty$ toward the boundary in any direction with $n_k>0$. Hence
any interior maximum has $\operatorname{rank}(\hat\rho_{\rm MLE})=d$. Asymptotic normality
of MLE requires the true parameter to be interior; boundary targets invalidate the delta method.
\end{proof}

\begin{compbox}
\textbf{Orthodox assumption:} Regularization (hedged MLE, MAP with Bures prior) enforces
positivity algebraically.

\textbf{RSVP replacement:} Positivity of the posterior emerges from $\Sfield>0$ without
regularization (Theorem~\ref{thm:rsvp_prior_well}).

\textbf{Orthodox pathology:} Regularization Jacobians bias estimates toward pure states.

\textbf{RSVP resolution:} The RSVP prior generates its own curvature through the entropy
field, with no Cholesky-induced bias.
\end{compbox}

% =============================================================
\part{The RSVP Framework}
% =============================================================

\chapter{Admissibility Geometry: Probability as Emergent Accessibility}
\label{ch:admissibility}

\begin{depbox}
This chapter introduces the central conceptual foundation of the RSVP framework. All
subsequent chapters in Parts~II--VI depend on Definitions~\ref{def:admissibility}--\ref{def:emergent_prob}
and Axiom~\ref{axiom:admissibility_primacy}.
\end{depbox}

\section{Trajectory Space and Representational Compression}

Let $\Traj$ be the space of all admissible world-histories and $\Rep$ a finite-dimensional
manifold on which statistical models operate. A \emph{projection} $\proj:\Traj\to\Rep$ maps
trajectory structure into observable quantities. For each $m\in\Rep$, the \emph{fiber}
$\proj^{-1}(m)$ is the set of all trajectories projecting to the same point. Classical
models are defined entirely on $\Rep$ and are silent about fiber structure.

\begin{definition}[Admissibility functional]
\label{def:admissibility}
An \emph{admissibility functional} $\Admiss:\Traj\times\mathbb{R}\to[0,\infty)$ measures
the degree to which trajectory $x$ remains dynamically and inferentially available at time $t$.
\end{definition}

\begin{definition}[Emergent probability]
\label{def:emergent_prob}
The \emph{emergent probability} of a representational event $B\subset\Rep$ is
\begin{equation}
  P(B,t) = \frac{\int_{\proj^{-1}(B)}\Admiss(x,t)\,dx}{\int_\Traj\Admiss(x',t)\,dx'}.
  \label{eq:emergent_prob}
\end{equation}
\end{definition}

\begin{axiom}[Admissibility primacy]
\label{axiom:admissibility_primacy}
The admissibility functional $\Admiss$ is the fundamental object of inference. Probability
distributions over $\Rep$ are derived quantities obtained by projecting $\Admiss$ via
equation~\eqref{eq:emergent_prob}.
\end{axiom}

\section{Entropy as Accessibility}

In the admissibility framework, entropy is the logarithm of accessible trajectory volume:
\begin{equation}
  \Sfield(x,t) = \log\Omega(x,t),
\end{equation}
where $\Omega(x,t)$ is the volume of admissible continuations from $x$ at time $t$. Low
entropy means strong constraint; high entropy means rich accessibility. This differs from
Shannon uncertainty, which is defined over predefined symbol sets.

\begin{proposition}[Shannon entropy as a limit]
\label{prop:shannon_limit}
If $\Admiss(x,t)=p(x)$ is time-independent and fibers are uniform, then the accessibility
entropy $\Sfield$ equals the Shannon entropy of the projected distribution $P(\cdot,t)$
up to a constant depending on fiber volume.
\end{proposition}

\begin{proof}
Under uniform fibers $|\proj^{-1}(m)|=C$, $\Omega(m,t)=Cp(m)$, so $\Sfield(m,t)=\log p(m)+\log C$.
Then $-\mathbb{E}[\Sfield]=H(P(\cdot,t))-\log C$, giving $H=\mathbb{E}[\Sfield]+\mathrm{const}$.
\end{proof}

\section{Why Bayesian Inference Fails at the Boundary of $\Rep$}

The projection $\proj$ is many-to-one. Inference fails at the boundary of $\Rep$ precisely
when the boundary corresponds to measure-zero trajectory fibers: the density matrix manifold
boundary (pure states) is the image of a measure-zero set of fibers in trajectory space,
yet maximum likelihood and Lindblad machineries assign finite probability to these boundary
regions by treating them as interior points of $\Rep$. This is the geometric explanation
for Proposition~\ref{prop:mle_boundary}.

\begin{compbox}
\textbf{Orthodox assumption:} The representational manifold $\Rep$ is a complete description;
all relevant information is in $m\in\Rep$.

\textbf{RSVP replacement:} $\Rep$ is a projection of $\Traj$; the fiber $\proj^{-1}(m)$
contains the lost causal history, alternative admissible paths, and trajectory distinctions
relevant to future inference.

\textbf{Orthodox pathology:} Fiber degeneracy causes observational aliasing,
boundary singularities, and representational fragility under model change.

\textbf{RSVP resolution:} The admissibility functional is defined on $\Traj$, making
fiber structure explicit and the cost of compression computable via the degeneracy ratio $\Gamma(m)$.
\end{compbox}

\section{Minimal Working Example: One-Dimensional RSVP System}

Before introducing the full RSVP machinery, we work through a minimal complete example
to anchor the formalism.

\begin{exbox}
\textbf{Binary Bayesian inference via 1D RSVP.} Let $\Theta=\{0,1\}$ encoded as
$\mathcal{M}=[-1,1]$, with $\Phi(\theta)=\log p(D\given\theta)$ initialized from a
binomial likelihood, $\Sfield(\theta)=(1-\theta^2)^{1/2}$ (high entropy near $\theta=0$,
low near the boundaries), and $\vfield(\theta)=(\alpha/\Sfield)\nabla\Phi$.

The RSVP prior is
$\pi_{\rm RSVP}(\theta)\propto(1-\theta^2)^{1/2}e^{\Phi(\theta)}$,
which is a semicircle-weighted likelihood. The posterior concentrates near the maximum of
$\Phi(\theta)$ but is moderated by the entropy field: if data strongly favor $\theta=+1$
(a boundary), the entropy field $\Sfield(\pm 1)=0$ suppresses the prior at the boundary,
preventing the overconfident boundary concentration of the MLE.

A three-event KES history $H_3=(e_1,e_2,e_3)$ with $e_i$ drawn from this system gives
an empirical frequency that converges to $\pi_{\rm RSVP}$ by Theorem~\ref{thm:kes_born}.
\end{exbox}

\chapter{The Relativistic Scalar-Vector Plenum: Field Equations}

\begin{depbox}
This chapter depends on Chapter~\ref{ch:admissibility} (admissibility geometry),
Chapter~2 (Riemannian manifolds and Fisher-Rao metric), and the Euler-Lagrange formalism.
\end{depbox}

\section{The RSVP Lagrangian}

\begin{definition}[RSVP Lagrangian density]
\begin{equation}
  \mathcal{L}_{\rm RSVP} = \tfrac{1}{2}|\nabla\Phi|^2
  - \tfrac{1}{2}\Sfield\,g(\vfield,\vfield)
  - V(\Phi,\Sfield)
  + \alpha\,\Phi\operatorname{div}(\vfield)
  - \beta\,\Sfield\log\Sfield.
  \label{eq:rsvp_lagrangian}
\end{equation}
\end{definition}

The term $-\beta\Sfield\log\Sfield$ drives $\Sfield$ dynamically toward maximum entropy
configurations, implementing the Jaynes principle as field evolution. The Euler-Lagrange
equations are:
\begin{align}
  \Delta\Phi &= \tfrac{\partial V}{\partial\Phi} - \alpha\operatorname{div}(\vfield),
  \label{eq:rsvp_phi}\\
  \Sfield\,\vfield^\flat &= \alpha\,d\Phi,
  \label{eq:rsvp_v}\\
  \partial_t\Sfield &= -\operatorname{div}(\Sfield\vfield)
  + \beta(1+\log\Sfield) + \tfrac{\partial V}{\partial\Sfield}.
  \label{eq:rsvp_S}
\end{align}

\begin{implbox}
\textbf{Numerical evolution.} Equations~\eqref{eq:rsvp_phi}--\eqref{eq:rsvp_S} are
discretized on a rectangular lattice with semi-Lagrangian advection for
$\operatorname{div}(\Sfield\vfield)$. Clamp $\Sfield\geq\epsilon>0$ at each step.
\begin{verbatim}
  Phi  += dt*(laplacian(Phi) - dV_dPhi + alpha*divergence(v))
  v     = (alpha/S)*gradient(Phi)
  S    += dt*(-divergence(S*v) + beta*(1+log(S)) + dV_dS)
  S     = clip(S, epsilon, None)
\end{verbatim}
\end{implbox}

\section{Well-Posedness and the RSVP Prior}

\begin{theorem}[Well-posedness]
\label{thm:rsvp_wellposed}
Let $(\mathcal{M},g)$ be compact without boundary, $\dim\leq 4$, and
$V=\frac{\lambda}{4}(\Phi^2-\mu^2)^2+\frac{\kappa}{2}\Phi\Sfield$.
For initial data $(\Phi_0,\vfield_0,\Sfield_0)\in H^2\times H^1\times H^1$ with
$\Sfield_0\geq s_0>0$, there exists a unique strong solution on $[0,T]$ with
$\Sfield\geq s_0 e^{-Ct}$.
\end{theorem}

\begin{proof}
Define energy $\mathcal{E}=\frac{1}{2}\|\nabla\Phi\|^2+\frac{1}{2}\|\Sfield^{1/2}\vfield\|^2
+\int V\,d\mathrm{vol}+\beta\int\Sfield\log\Sfield\,d\mathrm{vol}$.
Differentiating and integrating by parts on the compact manifold eliminates boundary terms.
Sobolev embedding $H^2\hookrightarrow L^\infty$ for $n\leq 4$ bounds the potential gradient.
Gronwall gives $\mathcal{E}(t)\leq\mathcal{E}(0)e^{C_1 t}$; local existence follows from
contraction mapping. The $\Sfield$ lower bound follows from the comparison principle applied
to~\eqref{eq:rsvp_S}.
\end{proof}

\begin{definition}[RSVP-generated prior]
At steady state, the RSVP prior over $\Theta\subset\mathcal{M}$ is
\begin{equation}
  \pi_{\rm RSVP}(\theta) = \frac{\Sfield(\theta)\,e^{\Phi(\theta)}}
  {\int_\Theta\Sfield(\theta')\,e^{\Phi(\theta')}\,d\mathrm{vol}_g}.
  \label{eq:rsvp_prior}
\end{equation}
\end{definition}

\begin{theorem}[RSVP prior positivity]
\label{thm:rsvp_prior_well}
Under Theorem~\ref{thm:rsvp_wellposed}, $\pi_{\rm RSVP}$ is strictly positive on $\Theta$
and defines a probability measure, resolving Proposition~\ref{prop:gaussian_failure}
at the dynamical level.
\end{theorem}

\begin{proof}
$\Sfield(\theta)e^{\Phi(\theta)}>0$ everywhere by $\Sfield>0$ and $e^\Phi>0$.
The denominator is a positive continuous integral over a compact set, hence finite and positive.
\end{proof}

\section{Correspondence Limits}

\begin{proposition}[Classical Bayesian limit]
If $\Sfield=\mathrm{const}$ and $\Phi=\log\mathcal{L}(\theta)$, then RSVP inference
reduces to ordinary Bayesian updating with likelihood $\mathcal{L}$ and uniform prior.
\end{proposition}

\begin{proposition}[Shannon entropy limit]
If $\Admiss(x,t)=p(x)$ (time-independent, uniform fibers), then the RSVP entropy field
$\Sfield$ equals the Shannon entropy of the projected distribution up to a fiber-volume constant
(Proposition~\ref{prop:shannon_limit}).
\end{proposition}

\begin{proposition}[Markovian limit]
The KES memory kernel $K_m$ (Definition~\ref{def:kes_memory}) reduces to the standard KES
map when $m=0$, giving Markovian event selection.
\end{proposition}

\begin{proposition}[Gaussian approximation limit]
Near a maximum $\theta^*$ of $\Phi$, Taylor expanding to second order gives
$\pi_{\rm RSVP}(\theta)\approx\Sfield(\theta^*)e^{\Phi(\theta^*)}\exp(-\frac{1}{2}(\theta-\theta^*)^T F(\theta^*)(\theta-\theta^*))$,
recovering a Gaussian posterior with Fisher information covariance.
\end{proposition}

\section{UV Finiteness of the Multimode Prior}

\begin{theorem}[UV finiteness]
\label{thm:uv_finite}
Let $\mathcal{M}=\prod_{k=1}^K[0,\infty)$ parameterize $K$ modes. With $\kappa_k\geq ck^\gamma$
for $c,\gamma>0$, the steady-state RSVP entropy field satisfies
$\int_\mathcal{M}\Sfield^*\,d\bar n<\infty$ without any ultraviolet cutoff.
\end{theorem}

\begin{proof}
The potential penalizes large $\bar n_k$ with strength $\kappa_k$, forcing
$\Sfield^*\leq C\prod_k e^{-\kappa_k\bar n_k}$ by the maximum principle.
Then $\int\Sfield^*\leq C\prod_k\kappa_k^{-1}$, which vanishes as $K\to\infty$
for $\kappa_k\geq ck^\gamma$.
\end{proof}

\chapter{RSVP Applied to Quantum Optical Inference}

\begin{depbox}
This chapter applies Theorems~\ref{thm:rsvp_wellposed}--\ref{thm:uv_finite} to concrete
photonic inference problems. Requires the Mandel $Q$ parameter and photon number operator
from Chapter~1.
\end{depbox}

\section{The Velocity Field as Natural Gradient}

\begin{theorem}[RSVP velocity as entropic natural gradient]
At steady state with $g=g^{\rm FR}$, the RSVP velocity field satisfies
$\vfield=(\alpha/\Sfield)(g^{\rm FR})^{-1}\nabla\Phi$:
the Fisher-Rao natural gradient of the likelihood, suppressed by entropy density in
high-uncertainty regions. Bayesian updating is geometrically implemented by the RSVP
velocity field. The RSVP framework does not add Bayesian inference to a field theory;
it shows that Bayesian inference \emph{is} a field theory with Lagrangian~\eqref{eq:rsvp_lagrangian}.
\end{theorem}

\begin{proof}
From~\eqref{eq:rsvp_v}, $\Sfield\vfield^\flat=\alpha\,d\Phi$. Raising with $(g^{\rm FR})^{-1}$
gives $\Sfield\vfield=\alpha\nabla_{g^{\rm FR}}\Phi$, so $\vfield=(\alpha/\Sfield)(g^{\rm FR})^{-1}\nabla\Phi$.
\end{proof}

\section{Mandel $Q$ Inference}

\begin{exbox}
\textbf{Sub-Poissonian light.} From $N=20$ counts with $\bar n=0.4$ and $s^2=0.1$,
the standard estimator gives $\hat Q=-0.75$. A Gaussian prior on $Q\in(-1,\infty)$ assigns
non-negligible mass to the unphysical region $Q<-1$. The RSVP prior with
$V\supset-\lambda\log(Q+1)$ gives $\Phi\to-\infty$ as $Q\to-1^+$, automatically enforcing
the constraint. The posterior mean is regularized toward $Q=0$ by $\Sfield$ weighting,
preventing over-interpretation of small-sample fluctuations.
\end{exbox}

\chapter{RSVP Photon Propagation and Entropic Phase}

\begin{depbox}
This chapter depends on Theorem~\ref{thm:rsvp_prior_well}, Section~2.1 (manifolds),
and the free photon propagator from quantum field theory.
\end{depbox}

\section{RSVP-Modified Propagator}

\begin{definition}[RSVP-modified propagator]
\begin{equation}
  \widetilde D_{\mu\nu}(x,y) = D_{\mu\nu}(x-y)
  \exp\!\left(-\tfrac{1}{2}\int_x^y\nabla\Phi\cdot d\ell
  -\tfrac{\alpha}{2}\int_x^y\Sfield\,ds\right).
  \label{eq:rsvp_propagator}
\end{equation}
\end{definition}

The potential factor $e^{-(\Phi(y)-\Phi(x))/2}$ is path-independent. The attenuation
factor $e^{-\frac{\alpha}{2}\int\Sfield\,ds}$ encodes decoherence without invoking the
environment explicitly: photons traversing high-entropy regions experience additional damping.

\begin{proposition}[RSVP entropic phase shift]
A photon of frequency $\omega_0$ propagating from $x_0$ to $x_1$ acquires
$\Delta\omega = \frac{\omega_0}{2}(\Phi(x_1)-\Phi(x_0))+O(\alpha\Sfield)$.
\end{proposition}

\begin{epistbox}
The identification of $\widetilde D_{\mu\nu}$ with the physical photon Green's function is
an ansatz. Testing the entropic phase shift requires a regime where $\nabla\Phi$ is
significant along the photon path, e.g., photonic crystals with engineered entropy profiles.
\end{epistbox}

\section{Optics as Admissibility Dynamics}

Standard optical phenomena acquire RSVP interpretations: interference is admissibility-bundle
compatibility; diffraction is the spread of admissible continuations after an aperture;
coherence $g^{(1)}(\tau)$ measures admissibility-bundle overlap over time $\tau$; a lens
is an admissibility operator imposing $\Phi(x)=-k|x|^2/(2f)$ that drives $\vfield$ toward
focal convergence.

\begin{exbox}
\textbf{Mach-Zehnder interferometer.} A photon entering path 1 or path 2 corresponds to
two trajectory bundles $x_1,x_2$ with $\proj(x_1)=\proj(x_2)$ at the final detector (same
outcome). Constructive interference occurs when the admissibility phases are compatible,
i.e., $\Phi(x_1)-\Phi(x_2)\in 2\pi\mathbb{Z}$. Destructive interference corresponds to
admissibility cancellation: the fiber $\proj^{-1}(\text{detector})$ has two contributions
that exactly cancel in the admissibility integral~\eqref{eq:emergent_prob}.
\end{exbox}

% =============================================================
\part{The Spherepop/KES Calculus}
% =============================================================

\chapter{The Spherepop Operators: An Irreversible Event Calculus}
\label{ch:spherepop}

\begin{depbox}
This chapter depends on Chapter~\ref{ch:admissibility} (admissibility geometry) and
the decoherence functional (Definition~\ref{def:decoherence_functional}).
\end{depbox}

\section{Event Primacy}

\begin{axiom}[Event primacy]
\label{axiom:event_primacy}
A state in the usual sense is an equivalence class of event histories. Two histories are
equivalent if and only if they are indistinguishable by all subsequent events.
\end{axiom}

Axiom~\ref{axiom:event_primacy} inverts the standard ordering: states are derived from
events, not the reverse. Irreversibility is primitive, not emergent.

\begin{definition}[Event algebra]
An \emph{event algebra} $(\mathcal{E},\leq,\sqcup,\bot)$ consists of events, causal
precedence, event composition, and null event. An event $e$ is \emph{irreversible} if no
$e^{-1}$ satisfies $e\sqcup e^{-1}=\bot$. All photon detection events are irreversible.
\end{definition}

\section{The Four Operators}

\begin{definition}[$\Pop$] Maps an event to its immediate causal consequence; models photon
detection as the creation of a new event, not a state update.\end{definition}

\begin{definition}[$\Bind$] Joins causally compatible events: $\Bind(e_1,e_2)=e_1\sqcup e_2$
when neither $e_1\leq e_2$ nor $e_2\leq e_1$; models coincidence measurement of entangled pairs.\end{definition}

\begin{definition}[$\Collapse$] Maps an event $e$ to the subspace
$\{|\psi\rangle:P_k|\psi\rangle\neq 0\}$ of states compatible with outcome $k$; identifies
compatible configurations without destroying pre-event information.\end{definition}

\begin{definition}[$\Refuse$] Filters events by a consistency condition $\mathcal{C}$:
$\Refuse(A,\mathcal{C})=\{e\in A:e\text{ is consistent with all }c\in\mathcal{C}\}$;
implements pointer basis selection.\end{definition}

\begin{theorem}[$\Refuse$ resolves the pointer basis problem]
\label{thm:refuse_pointer}
With $\mathcal{C}=\{\text{robustness under tracing out }\mathcal{H}_E\}$,
$\Refuse(\mathcal{E}_S,\mathcal{C})$ selects precisely the Quantum Darwinism pointer basis:
system states that leave redundant records in multiple environmental fragments.
\end{theorem}

\begin{proof}
A state is in the pointer basis iff $I(S:E_\alpha)\approx H(S)$ for many disjoint
environmental fragments $\mathcal{H}_{E_\alpha}$. The $\Refuse$ consistency condition
requires any event associated with the state to be reproducible from multiple fragments,
which is precisely the Quantum Darwinism redundancy condition.
\end{proof}

\begin{proposition}[Non-commutativity of $\Pop$ and $\Bind$]
$\Pop(\Bind(e_1,e_2))\neq\Bind(\Pop(e_1),\Pop(e_2))$ in general.
\end{proposition}

\begin{proof}
For photons from parametric down-conversion, $\Pop(\Bind(e_1,e_2))$ yields HOM-correlated
outcomes; $\Bind(\Pop(e_1),\Pop(e_2))$ yields no HOM interference. The measurement statistics
differ experimentally, confirming the inequality.
\end{proof}

\chapter{KES: Kinetic-Event Synthesis}
\label{ch:kes}

\begin{depbox}
This chapter depends on the Spherepop operators (Chapter~\ref{ch:spherepop}), the RSVP
prior~\eqref{eq:rsvp_prior}, and the admissibility functional (Definition~\ref{def:admissibility}).
\end{depbox}

\section{The KES Map}

\begin{definition}[World-state and history state]
$\Omega_t=(\mathcal{X}_t,w_t)$ with $w_t:\mathcal{X}_t\to[0,1]$ normalized is the
posterior-weighted set of consistent configurations at time $t$.
$H_t=(e_1,\ldots,e_m)$ is an irreversible causal sequence.
\end{definition}

\begin{definition}[KES map]
\label{def:kes}
The KES map $\mathrm{KES}:\Omega_t\to H_{t+1}$ selects $e^*$ maximizing
\begin{equation}
  \mathcal{F}(e) = w_t(\Collapse(e))\cdot\Sfield(\Collapse(e))\cdot e^{\Phi(\Collapse(e))},
  \label{eq:kes_functional}
\end{equation}
and sets $H_{t+1}=H_t\sqcup(e^*)$.
\end{definition}

\begin{theorem}[KES recovers Born statistics]
\label{thm:kes_born}
When $\Sfield=\mathrm{const}$ and $\Phi=\log p$, the long-run frequency of KES-selected
outcomes converges almost surely to $p(e_k)=\operatorname{tr}(\rho E_k)$.
\end{theorem}

\begin{proof}
$\mathcal{F}(e_k)=S_0 w_t(\Collapse(e_k))p(e_k)$; for uniform $w_t$, KES samples from
$(p(e_k))_k$. The strong law of large numbers gives almost-sure convergence to
$\operatorname{tr}(\rho E_k)$.
\end{proof}

\section{KES Update and No-Collapse Interpretation}

Post-selection updates $w_t$ by Bayesian conditioning on $e^*$:
\begin{equation}
  w_{t+1}(x) = \frac{w_t(x)\cdot\mathbf{1}[x\in\Collapse(e^*)]}
  {\sum_{x'\in\Collapse(e^*)}w_t(x')}.
\end{equation}

No Hilbert-space collapse occurs. The apparent collapse is an artifact of conflating history
states with world-states.

\section{KES Entropy and the Second Law}

\begin{definition}[KES entropy]
\begin{equation}
  \mathcal{S}_{\rm KES}(t) = -\sum_x w_t(x)\log w_t(x) + |H_t|\log|\mathcal{E}|.
\end{equation}
\end{definition}

\begin{theorem}[KES second law]
$\mathbb{E}[\mathcal{S}_{\rm KES}(t+1)]\geq\mathcal{S}_{\rm KES}(t)$, with equality iff
$w_t$ is a point mass.
\end{theorem}

\begin{proof}
The history term increases by $\log|\mathcal{E}|>0$ per step. The Shannon term decreases
by at most $H(w_t)$ by the data processing inequality. The net change is non-negative for
$|\mathcal{E}|\geq|\mathcal{X}_t|$.
\end{proof}

\section{Non-Markovian Memory}

\begin{definition}[KES memory kernel of depth $m$]
\label{def:kes_memory}
\begin{equation}
  K_m(e_{t-m},\ldots,e_{t-1}) = \int_{\mathcal{X}_t}
  w_t(x\given H_{t-m:t})\,\mathcal{F}(e^\star(x))\,dx.
  \label{eq:kes_memory}
\end{equation}
\end{definition}

For $m=0$, the kernel is Markovian. For $m>0$, past events influence event selection,
modeling the distinguishability revivals of Theorem~\ref{thm:non_markov_entropy}.

\begin{implbox}
\textbf{KES memory as transformer attention.} The memory kernel $K_m$ is approximated by
attention over the last $m$ events encoded as feature vectors. This maps KES dynamics
directly onto transformer architecture: attention heads may be learning approximations to
admissibility kernels.
\end{implbox}

% =============================================================
\part{Synthesis and Connections}
% =============================================================

\chapter{Unification: RSVP Fields as KES Selection Geometry}

\begin{depbox}
This chapter synthesizes all preceding material. Requires Theorems~\ref{thm:rsvp_wellposed},
\ref{thm:rsvp_prior_well}, \ref{thm:kes_born}, and the Spherepop operators.
\end{depbox}

\section{The Unification Theorem}

\begin{theorem}[RSVP-KES unification]
\label{thm:unification}
Let $(\mathcal{M},g,\rsvp)$ be a compact RSVP configuration and define the KES functional
by~\eqref{eq:kes_functional}. Then: (1) KES-selected events form a Markov chain with
kernel $K(e,e')=\mathcal{F}(e')/Z(e)$; (2) the stationary distribution is $\pi_{\rm RSVP}$;
(3) convergence is exponential with rate given by the spectral gap of the RSVP Laplacian.
\end{theorem}

\begin{proof}
(1) $\mathcal{F}(e')>0$ and $Z(e)<\infty$ on a compact manifold give a well-defined kernel.
(2) At steady state $\mathcal{F}(\theta)\propto\Sfield e^\Phi$ and $Z$ is constant; detailed
balance confirms $\pi=\pi_{\rm RSVP}$. (3) The Poincar\'{e} inequality on a compact
Riemannian manifold gives a positive spectral gap for $\Delta$, which controls the Markov
chain mixing time.
\end{proof}

\chapter{Comparisons: QBism, Darwinism, Entropic Dynamics, Free Energy}

\begin{depbox}
This chapter compares RSVP-KES to four existing programs. Requires Theorems~\ref{thm:kes_born},
\ref{thm:refuse_pointer}, and the ELBO~\eqref{eq:elbo}.
\end{depbox}

For each framework, we state what RSVP-KES reproduces, what it modifies, and what
predictions differ.

\section{QBism}

\textbf{Reproduced:} Epistemic status of $\rho$; Born-rule consistency in the flat limit
(Theorem~\ref{thm:kes_born}).
\textbf{Modified:} The consistency norm is replaced by the KES functional dynamics. The
SIC-POVM correction term acquires a field-theoretic account via $\Sfield$ geometry.
\textbf{Differing predictions:} In non-trivial RSVP backgrounds, effective probabilities
deviate from Born by $O(\alpha\Sfield)$. This is absent in QBism, which provides no
field dynamics.

\section{Quantum Darwinism}

\textbf{Reproduced:} Theorem~\ref{thm:refuse_pointer} shows $\Refuse(\mathcal{E}_S,\mathcal{C}_{\rm robust})$
selects precisely the pointer basis.
\textbf{Modified:} The RSVP entropy field $\Sfield$ provides a dynamical measure of
redundancy; low $\Sfield$ corresponds to high environmental recording.
\textbf{Differing predictions:} A specific functional form for the quantum-to-classical
transition as a function of $\Sfield$ along the system-environment interaction.

\section{Entropic Dynamics}

\textbf{Reproduced:} The term $-\beta\Sfield\log\Sfield$ in the RSVP Lagrangian
implements maximum entropy dynamics, paralleling Caticha's entropic update.
\textbf{Modified:} Entropic Dynamics does not couple entropy to $\Phi$ or $\vfield$.
The RSVP velocity field has no analog in Entropic Dynamics.
\textbf{New structure:} The KES memory kernel~\eqref{eq:kes_memory} handles non-Markovian
dynamics that Entropic Dynamics cannot model.

\section{Friston's Free Energy Principle}

\textbf{Reproduced:} ELBO maximization~\eqref{eq:elbo} corresponds to minimizing
$\mathrm{KL}(q_\phi\|\pi_{\rm RSVP})$. Active inference corresponds to KES event selection.
\textbf{Modified:} The FEP treats the generative model $p$ as fixed. In RSVP-KES, the
prior $\pi_{\rm RSVP}$ is dynamically generated; evidence updates both the posterior and
the prior geometry.
\textbf{Formal connection:} The FEP is the special case of RSVP-KES where the fields are
static (no field dynamics, only posterior updates).

\chapter{Experimental Predictions and Falsifiability}

\begin{depbox}
This chapter requires the RSVP propagator~\eqref{eq:rsvp_propagator}, the KES functional,
and the memory kernel~\eqref{eq:kes_memory}. Predictions 1--3 depend on ansätze stated explicitly.
\end{depbox}

\begin{description}
\item[\textbf{Prediction 1}] \textit{Modified photon detection statistics.} In a photonic
  system where $\Sfield$ is engineered to be non-uniform, KES predicts
  $\tilde p_k\propto\Sfield_k\operatorname{tr}(\rho E_k)$, deviating from Born by an
  $O(\Sfield)$ factor. \textit{Ansatz:} identification of $\mathcal{F}$ with physical
  detection weight.

\item[\textbf{Prediction 2}] \textit{Entropic phase shift.} Photons propagating through a
  medium with $\nabla\Phi\neq 0$ acquire $\Delta\omega=\frac{\omega_0}{2}\nabla\Phi\cdot\Delta x$,
  distinguishable from dispersive shifts by its $\nabla\Phi$ dependence.

\item[\textbf{Prediction 3}] \textit{Controlled non-Markovian entropy revivals.} Systems
  with $m>0$ KES memory exhibit distinguishability revivals at times determined by $K_m$
  without free parameters once $\Sfield$ and $\Phi$ are calibrated.

\item[\textbf{Prediction 4}] \textit{UV-finite multimode prior.} The RSVP prior is
  normalizable without UV cutoff (Theorem~\ref{thm:uv_finite}), predicting specific
  mode-number dependence of prior probability testable against Jeffreys and Bures priors.

\item[\textbf{Prediction 5}] \textit{Pointer basis selection by $\Refuse$.} In quantum eraser
  experiments, $\Refuse$ predicts which basis is selected with a quantitative relationship
  between $\Sfield$ and the redundancy fraction, agreeing with Quantum Darwinism but
  providing a specific dynamical mechanism.
\end{description}

\begin{epistbox}
Predictions 1 and 2 require experimentally accessing RSVP field quantities independently.
Current proposals: (a) inverse RSVP field equations from posterior data; (b) deliberate
engineering of $\Sfield$ via structured reservoir coupling; (c) estimating $\Phi$ from the
natural gradient structure of the experimental likelihood. These are research directions,
not established methods.
\end{epistbox}

% =============================================================
\part{Extensions}
% =============================================================

\chapter{Projection Failure and Representational Collapse}

\begin{depbox}
This chapter depends on Chapter~\ref{ch:admissibility} (admissibility geometry and the
projection $\proj:\Traj\to\Rep$), Proposition~\ref{prop:mle_boundary}, and
Theorem~\ref{thm:non_markov_entropy}.
\end{depbox}

\section{Fiber Degeneracy and Observational Aliasing}

Define the \emph{degeneracy ratio}
\begin{equation}
  \Gamma(m) = \frac{\mu(\proj^{-1}(m))}{\int_\Rep\mu(\proj^{-1}(m'))\,dm'},
\end{equation}
measuring the relative volume of the fiber over $m\in\Rep$. Large $\Gamma(m)$ means many
distinct trajectories are aliased to the same representational point: the inference system
cannot distinguish them, and any prediction that depends on their difference is compromised.

\begin{proposition}[Degeneracy implies inference failure]
If $\Gamma(m)>\Gamma(m')$ and future predictions depend on trajectories in
$\proj^{-1}(m)\setminus\proj^{-1}(m')$, then any model operating on $\Rep$ cannot achieve
the prediction accuracy of a model operating on $\Traj$.
\end{proposition}

\begin{proof}
By definition, the fiber $\proj^{-1}(m)$ contains trajectories that are indistinguishable
from $m$ alone. If these trajectories lead to different future outcomes, then any model
that uses only $m$ cannot distinguish the outcomes and must average over the fiber,
producing prediction error equal to the variance of outcomes within the fiber.
\end{proof}

\section{Boundary Singularities as Measure-Zero Fibers}

The boundary of $\Rep$ (e.g., pure states in quantum tomography) corresponds to
measure-zero fibers in $\Traj$: only a single trajectory (or a lower-dimensional
submanifold) projects to each boundary point. The MLE boundary concentration
(Proposition~\ref{prop:mle_boundary}) is therefore a symptom of a deeper pathology:
the inference system is assigning finite probability to representational regions
whose preimage in trajectory space has measure zero.

\section{Projection Curvature and Berry Phase}

The curvature of the fiber bundle $\proj:\Traj\to\Rep$ is measured by the \emph{connection}
induced by the RSVP vector field $\vfield$ on the horizontal distribution
$H_x\Traj = \ker(d\proj_x)^\perp$. Non-trivial curvature---the holonomy of this
connection around closed loops in $\Rep$---corresponds to the Berry phase in quantum
mechanics: a photon transported around a loop in parameter space acquires a geometric
phase proportional to the fiber bundle curvature.

\begin{conjecture}
The Berry phase of a quantum optical system is equal to the holonomy of the RSVP
admissibility connection $\nabla^{\rm RSVP}$ around the corresponding loop in
the statistical manifold $\Rep$.
\end{conjecture}

\begin{implbox}
\textbf{Sparse tensor representation of high-dimensional fibers.} For $K$-mode systems,
$\Traj$ has dimension $O(4^K)$ while $\Rep$ has dimension $O(K)$. The fiber
$\proj^{-1}(m)$ is an $O(4^K-K)$-dimensional manifold. Sparse tensor formats (Tucker
decomposition, tensor train) can represent the admissibility measure on fibers with
$O(Kr^2)$ parameters for bond dimension $r$, enabling scalable fiber-aware inference.
\end{implbox}

\chapter{Information Geometry Beyond Fisher-Rao}

\begin{depbox}
This chapter extends Chapter~2 (Fisher-Rao geometry) using the RSVP entropy field from
Chapter~\ref{ch:admissibility} and the field equations~\eqref{eq:rsvp_phi}--\eqref{eq:rsvp_S}.
\end{depbox}

\section{Limitations of Fisher-Rao Geometry}

The Fisher-Rao metric assumes: fixed probability spaces, predefined observables, and static
admissibility. It is a metric on \emph{representational} manifold $\Rep$, not on the
trajectory space $\Traj$. It captures the geometry of probability distributions but not
the geometry of the fibers that generate them.

\section{The RSVP Metric}

\begin{definition}[RSVP metric]
The RSVP metric on $\mathcal{M}$ is
\begin{equation}
  g^{\rm RSVP}_{ij} = g^{\rm FR}_{ij} + \lambda\,\partial_i\Sfield\,\partial_j\Sfield,
  \label{eq:rsvp_metric}
\end{equation}
extending the Fisher-Rao metric by a term encoding entropy gradients.
\end{definition}

Entropy gradients contribute to the metric because regions where $\Sfield$ changes rapidly
are informationally distinct: inference is harder across entropy-gradient boundaries. The
RSVP metric penalizes such boundaries geometrically.

\begin{proposition}[RSVP metric is positive definite]
For $\lambda\geq 0$, $g^{\rm RSVP}$ is positive semidefinite; for $\lambda>0$ and
$\Sfield$ non-constant, it is positive definite.
\end{proposition}

\begin{proof}
$g^{\rm FR}$ is positive semidefinite by construction. The second term
$\lambda\,\partial_i\Sfield\,\partial_j\Sfield$ is a rank-one positive semidefinite correction.
For $\lambda>0$ and $\nabla\Sfield\neq 0$, the correction is positive definite in the
direction of $\nabla\Sfield$, and $g^{\rm FR}>0$ in the orthogonal directions.
\end{proof}

\section{Admissibility Geodesics}

Geodesics in the RSVP metric~\eqref{eq:rsvp_metric} generalize Fisher-Rao geodesics
(which correspond to the most informatively efficient paths between distributions) by
penalizing paths that cross entropy-gradient barriers. These geodesics correspond to:
\begin{itemize}
\item Bayesian natural gradient paths (in the Fisher-Rao component),
\item Entropy-avoiding inference trajectories (in the $\lambda\,\partial\Sfield\,\partial\Sfield$ component),
\item Wasserstein gradient flows when the metric is the $L^2$ Wasserstein metric on the space of probability measures.
\end{itemize}

\begin{exbox}
\textbf{Geodesic on the Bloch sphere.} For a single qubit, the Fisher-Rao metric on the
Bloch sphere is the round metric $g^{\rm FR}=d\theta^2+\sin^2\theta\,d\phi^2$. The RSVP
metric adds a term $\lambda(\partial_\theta\Sfield)^2 d\theta^2+\lambda(\partial_\phi\Sfield)^2 d\phi^2$.
If $\Sfield(\theta)=\sin\theta$ (maximum entropy at the equator, zero at poles), geodesics
are deflected toward the equatorial plane, reflecting the higher admissibility near $\theta=\pi/2$
and the lower admissibility (zero entropy) at the poles (pure states).
\end{exbox}

\chapter{Non-Markovian Dynamics and Recursive Inference}

\begin{depbox}
This chapter depends on Theorem~\ref{thm:non_markov_entropy}, the KES memory kernel
(Definition~\ref{def:kes_memory}), and Section~2.3 (entropy and variational principles).
\end{depbox}

\section{Inference With Historical Retention}

The KES memory kernel $K_m$ from Definition~\ref{def:kes_memory} defines recursive Bayesian
accessibility: the admissibility at time $t$ is modulated by the event history $H_{t-m:t}$,
\begin{equation}
  P_t(B) = \frac{\int_{\proj^{-1}(B)}\Admiss(x,t,H_t)\,dx}
  {\int_\Traj\Admiss(x',t,H_t)\,dx'}.
\end{equation}

This generalizes equation~\eqref{eq:emergent_prob} to history-dependent admissibility,
enabling systems that fail under Markovian approximation---photonic memory, biological
adaptation, recursive learning---to be modeled within the RSVP framework.

\section{Admissibility Continuity Equation}

\begin{theorem}[Admissibility conservation]
If $\operatorname{div}(\vfield)=0$, the total admissibility measure
$\mathfrak{A}(t)=\int_\mathcal{M}\Admiss(x,t)\,dx$ is conserved:
$\partial_t\mathfrak{A}=0$.
\end{theorem}

\begin{proof}
The admissibility continuity equation is
$\partial_t\Admiss+\operatorname{div}(\Admiss\vfield)=0$.
Integrating over $\mathcal{M}$ and applying the divergence theorem gives
$\partial_t\mathfrak{A}=-\int_{\partial\mathcal{M}}\Admiss\vfield\cdot\hat n\,dA$.
On a compact manifold without boundary, or with $\Admiss\vfield\cdot\hat n=0$ on $\partial\mathcal{M}$,
the boundary term vanishes.
\end{proof}

\chapter{Photonics as Constraint Propagation}
\label{ch:photonics_constraint}

\begin{depbox}
This chapter applies the RSVP propagator~\eqref{eq:rsvp_propagator} and the admissibility
framework (Chapter~\ref{ch:admissibility}) to optical systems from Chapter~1.
\end{depbox}

\section{Optical Systems as Admissibility Transformers}

An optical system (beam splitter, cavity, lens, fiber) is an operator on the admissibility
field that transforms $\rsvp$ while preserving certain conserved quantities.

\textbf{Beam splitters} partition the trajectory admissibility: a 50:50 beam splitter
performs $\Admiss_{\rm out}\propto\Admiss_{\rm in}^{(1)}+\Admiss_{\rm in}^{(2)}$
with interference structure in the admissibility overlap.

\textbf{Optical cavities} create low-entropy trajectory basins: resonant frequencies
correspond to stable fixed points of the RSVP velocity field $\vfield$, where
$\vfield\cdot\nabla\Phi=0$ (no net flow) and $\Sfield$ is locally minimized (high constraint).

\textbf{Wigner negativity} corresponds to inaccessible classical fibers: the negative
regions of $W(\alpha)$ are those where classical admissibility $\Admiss_{\rm cl}$ is
suppressed but quantum admissibility $\Admiss_{\rm qu}$ is not, reflecting the inaccessibility
of the classical trajectory structure in those phase-space regions.

\begin{exbox}
\textbf{Double-slit decoherence.} In a double-slit experiment, two trajectory bundles
$x_{\rm slit\,1}$ and $x_{\rm slit\,2}$ arrive at the same screen point. Their admissibility
overlap $\int\Admiss(x_{\rm slit\,1},t)\Admiss(x_{\rm slit\,2},t)\,dt$ determines the
fringe visibility. As the entropy field $\Sfield$ between the slits and the screen increases
(due to environmental coupling), the overlap decreases. The decoherence rate is
$\Gamma_D = \partial_t D(\rho_1,\rho_2) \propto \alpha\int\Sfield\,ds$, which is the
entropy-attenuation term in the RSVP propagator~\eqref{eq:rsvp_propagator}.
\end{exbox}

\chapter{Adversarial Robustness of Photonic Neural Networks: Non-Ideality as Defense}
\label{ch:adversarial}

\begin{depbox}
This chapter depends on Chapter~\ref{ch:admissibility} (admissibility geometry and the
projection $\proj:\Traj\to\Rep$), Chapter~\ref{ch:photonics_constraint} (photonics as
constraint propagation), and the KES sensitivity functional~\eqref{eq:kes_functional}.
It draws on the work of Lu et al.\ (2024), which demonstrated that hardware
non-idealities in photonic analog neural networks can function as intrinsic adversarial
defenders rather than merely as sources of degradation.
\end{depbox}

\section{Hardware Non-Ideality as an Admissibility Constraint}

Analog optical neural networks (ONNs) are photonic hardware systems in which neural
network weights are encoded as light amplitudes using digital-to-analog conversion (DAC)
modules, propagated through photonic tensor cores, and accumulated at photodetectors. The
standard design concern is that hardware non-idealities---low-precision control, optical
noise, crosstalk, fabrication variation---degrade inference accuracy relative to the
ideal digital computation. The conventional response is to suppress these non-idealities
through hardware-software co-design.

The RSVP framework suggests a different interpretation. Hardware non-idealities are
constraints on the admissibility field of the photonic system: they reduce the set of
reachable states $\mathcal{X}_t$ and alter the entropy field $\Sfield$ by restricting
accessible trajectory continuations. From the RSVP perspective, a system with low
precision has a coarser admissibility resolution---fewer distinct states are reachable---
which is a reduction in $\Sfield$ locally. This reduced accessibility is conventionally
undesirable for accuracy, but it simultaneously restricts the space available to an
adversary performing a weight attack.

Lu et al.\ (2024) made this precise for a specific class of attacks: gradient-based
bit-flip attacks (BFA) on stored neural network weights. Their central observation is
that low-bit quantization, sparse (pruned) weight representations, and the unary encoding
natively used by optical DAC hardware each reduce the \emph{weight sensitivity} to
bit-flip perturbations, and that this reduction has an adversarial protection effect that
can be harnessed deliberately. We formalize this observation within the RSVP framework
and extend it to a general principle.

\begin{definition}[Weight sensitivity in the ONN context]
\label{def:weight_sensitivity}
Let $\mathcal{L}(W, \mathcal{D})$ be the loss of an ONN with weights $W$ on dataset
$\mathcal{D}$. The \emph{bit-flip sensitivity} of weight $W_i$ is the Taylor expansion
of the loss change induced by flipping the most significant bit (MSB) of $W_i$:
\begin{equation}
  S_i = \mathcal{L}(W \oplus \Delta W_i^{\rm MSB}) - \mathcal{L}(W)
  \approx \nabla_{W_i}\mathcal{L}\cdot\Delta W_i^{\rm MSB}
  + \tfrac{1}{2}\nabla^2_{W_i}\mathcal{L}\cdot(\Delta W_i^{\rm MSB})^2,
  \label{eq:onn_sensitivity}
\end{equation}
where $\Delta W_i^{\rm MSB}$ is the weight perturbation caused by an MSB flip.
\end{definition}

In binary-coded decimal (BCD) representation with $b$ bits, an MSB flip changes the
weight by approximately half the weight range: $|\Delta W_i^{\rm MSB}| \approx 2^{b-1}$.
This is a large perturbation, and $S_i$ is correspondingly large for most weights. The
adversarial attack selects the weight with the largest $|S_i|$ and flips it, repeating
this process until all of a Hamming distance budget $HD$ is exhausted.

\section{Unary Encoding as Admissibility Restriction}

The key hardware observation is that optical DACs in photonic accelerators natively
encode weights in \emph{unary} format rather than BCD. In a unary-encoded $b$-bit weight,
the value $w$ is represented as $w$ leading ones followed by $(2^b - 1 - w)$ trailing
zeros: $(w)_B = (\underbrace{1,\ldots,1}_{w},\underbrace{0,\ldots,0}_{2^b-1-w})_U$.

\begin{proposition}[Unary encoding minimizes MSB sensitivity]
\label{prop:unary_sensitivity}
In the unary representation, every bit is a least-significant bit (LSB) in the sense
that flipping any single bit changes the weight value by exactly $\pm 1$ (one count unit).
In BCD representation, flipping the MSB changes the weight by $2^{b-1}$. Therefore,
$|\Delta W_i^{\rm MSB}|_{\rm unary} \ll |\Delta W_i^{\rm MSB}|_{\rm BCD}$ for $b \geq 2$,
and the sensitivity $S_i$ is correspondingly reduced by a factor of order $2^{b-1}$.
\end{proposition}

\begin{proof}
In a $b$-bit BCD representation with 2's complement, the MSB contributes value $-2^{b-1}$
(sign bit). An MSB flip changes the weight by $\pm 2^{b-1}$. In unary representation of
a $(2^b-1)$-bit string, each bit contributes value $+1$ (one count unit). Any single-bit
flip changes the encoded count by exactly $\pm 1$. Since the weight range is the same
in both representations, the ratio of MSB-flip perturbation sizes is $2^{b-1}:1$.
For $b=8$, this is a $128\times$ reduction in $|\Delta W^{\rm MSB}|$ and therefore
in the first-order sensitivity contribution.
\end{proof}

In RSVP terms, the unary encoding imposes a constraint on the admissibility field of the
weight space: the topology of the weight manifold $\mathcal{M}_W$ changes from one where
large jumps are accessible via single bit-flips (BCD) to one where only unit steps are
reachable (unary). This is a reduction in the local admissibility radius of the adversary's
action space. The entropy field $\Sfield$ of the weight manifold, which measures the volume
of accessible perturbations at each weight, is dramatically lower under unary encoding.

\begin{compbox}
\textbf{Orthodox assumption:} Hardware non-idealities (low precision, sparse encoding,
unary representation) are constraints that degrade accuracy and should be minimized.

\textbf{RSVP replacement:} Non-idealities are admissibility constraints on the weight
manifold $\mathcal{M}_W$. They reduce $\Sfield$ locally, restricting the trajectory space
available to an adversary. The reduction in adversarial accessibility is the dual of the
reduction in representational capacity.

\textbf{Orthodox pathology:} The adversary's action space (all possible MSB flips) is
large in BCD format; small Hamming distance attacks can cause catastrophic accuracy drops.

\textbf{RSVP resolution:} Unary encoding contracts the admissibility radius of the
adversary's actions to $O(1)$ rather than $O(2^{b-1})$, making even large Hamming
distance budgets ($HD=100$) nearly ineffective.
\end{compbox}

\section{Memory-Efficient Truncated Complementary Unary Encoding}

The original unary representation requires $2^b - 1$ bits to encode a $b$-bit weight,
a factor of $2^{b-1}-1/b$ memory overhead relative to BCD. This is prohibitive for
large networks. Lu et al.\ (2024) propose a memory-efficient variant, the
\emph{truncated complementary unary} (TCU) encoding, which exploits the Gaussian-like
weight distribution of trained neural networks.

The key observation is that the trailing zeros of small positive values and the leading
ones of small negative values are redundant for storage: they carry no information beyond
alignment. Truncating them reduces the required bitwidth from $2^b-1$ to $\hat b$, where
$\hat b$ is chosen to cover the weight's magnitude. For negative values, the complementary
encoding stores trailing zeros (counting from the right) rather than leading ones,
achieving the same compression for the negative half of the weight distribution.

\begin{proposition}[TCU memory overhead]
For a weight $W_i$ with $|W_i| = k$ (in integer count units, $0 \leq k \leq 2^{b-1}$),
the TCU representation requires
\begin{equation}
  \hat b_i = \lceil\log_2\min(2^b - |W_i|, |W_i|)\rceil + 1
  \label{eq:tcu_bits}
\end{equation}
bits, compared to $2^b - 1$ bits for the full unary representation. For weights
concentrated near zero (as in trained networks), $\hat b_i \ll 2^b - 1$ for most $i$,
giving memory overhead reduction of $6$--$11\times$ relative to full unary encoding.
\end{proposition}

The TCU encoding creates an exponentially-spaced partition of the weight range: weights
of magnitude $k$ require $O(\log_2 k)$ bits. This is analogous to logarithmic
quantization, but motivated by the unary admissibility constraint rather than by
rate-distortion considerations.

In RSVP terms, the TCU encoding defines a non-uniform admissibility metric on the weight
manifold: small weights have a fine-grained admissibility structure (many distinguishable
states in a small neighborhood) while large weights have a coarser structure. The
entropy field $\Sfield$ is correspondingly non-uniform, higher for large weights and
lower for small weights. Since empirically the most attack-sensitive weights tend to have
small magnitudes, the TCU encoding concentrates the low-$\Sfield$ region precisely
where adversarial vulnerability is highest.

\begin{exbox}
\textbf{TCU memory overhead for 8-bit VGG-8.} For a network where weights are stored
in 8-bit BCD, the full unary representation requires $2^8-1=255$ bits per weight,
a $255/8\approx 32\times$ memory overhead. With TCU encoding, protecting $0.2\%$ of
weights (the most sensitive ones) requires only $1.07\%$ additional memory overhead ---
a $30\times$ reduction. Lu et al.\ (2024) report that this $0.2\%$ protection rate,
combined with post-attack weight locking (Section~\ref{sec:weight_locking}), achieves
$86.73\%$ mean post-attack accuracy on VGG-8/CIFAR-10, compared to $32.89\%$ without
any defense, at a total memory overhead below $3\%$.
\end{exbox}

\section{Post-Attack Recovery via Sensitivity-Aware Weight Locking}
\label{sec:weight_locking}

Pre-attack unary protection covers only a small fraction $\alpha$ of weights. An adversary
with Hamming distance budget $HD$ will eventually attack unprotected weights.
Post-attack recovery requires: (1) detection of which weights were flipped, and
(2) correction of the flipped values.

Detection uses group-wise MSB checksum verification: weights are partitioned into groups
of size $G$; a checksum mismatch flags the entire group as potentially compromised.
The choice of $G$ trades detection precision (smaller $G$ localizes attacks more
precisely) against memory overhead (smaller $G$ requires more checksum storage).

\begin{definition}[Weight locking]
\label{def:weight_locking}
Given $K$ pre-computed \emph{centroids} $\{W_k\}_{k=1}^K$ for a layer's weights,
\emph{weight locking} replaces all weights in a detected victim group with their
nearest centroid. Centroids are computed pre-deployment via a sensitivity-aware
$K$-means variant in which the clustering distance is the locking-induced accuracy loss
(equation~\eqref{eq:onn_sensitivity}) rather than Euclidean distance:
\begin{equation}
  d_{in} = \nabla_{W_i}\mathcal{L}\cdot(W_i - \widetilde{W}_n)
  + \tfrac{1}{2}\nabla^2_{W_i}\mathcal{L}\cdot(W_i - \widetilde{W}_n)^2,
  \label{eq:locking_distance}
\end{equation}
where $\widetilde{W}_n$ is the centroid of the $n$-th detection group.
\end{definition}

Weight locking generalizes the simpler pruning-based recovery method, which forces
detected victim weights to zero. Pruning to zero is a special case of weight locking
with $K=1$ and centroid $W_1=0$. The locking approach allows $K>1$ centroids that
better approximate the pre-attack weight distribution, recovering higher accuracy
at comparable or lower memory cost.

The memory overhead of the locking scheme is
\begin{equation}
  m_L = \frac{|W|(\log_2 K + 2)}{G \cdot b|W|} \quad (G > 1),
\end{equation}
where the $\log_2 K$ term stores the cluster ID per weight and the $2/G$ term stores
the checksum. The accuracy-constrained optimization over $(G, K)$ is solved offline
by a greedy search that initializes at $(G=512, K=1)$ and halves $G$ while doubling
$K$ until the accuracy drop threshold $\eta$ is satisfied.

\section{RSVP Interpretation: Non-Ideality as Admissibility Engineering}

The full defense framework of Lu et al.\ (2024) --- TCU pre-attack encoding combined
with sensitivity-aware post-attack weight locking --- is, in RSVP terms, a deliberate
engineering of the admissibility field of the weight manifold $\mathcal{M}_W$.

The TCU encoding restricts the admissibility radius of the adversary's action space
(Proposition~\ref{prop:unary_sensitivity}): it lowers $\Sfield$ in the neighborhood of
vulnerable weights, making large-perturbation trajectories inaccessible to the attacker.
The weight locking pre-assigns centroid attractors in $\mathcal{M}_W$: after an attack,
the KES-like recovery map sends detected victim weights to their pre-designated centroid,
analogous to a $\Refuse$ operation that filters post-attack configurations to those
consistent with the pre-deployment constraint set.

\begin{proposition}[Weight locking as a Refuse operation]
\label{prop:locking_as_refuse}
Let $\mathcal{C}_{\rm lock} = \{W_k\}_{k=1}^K$ be the pre-deployed centroid set. Then
weight locking is equivalent to applying $\Refuse(\mathcal{E}_W, \mathcal{C}_{\rm lock})$:
it filters the post-attack weight configuration to the nearest element of
$\mathcal{C}_{\rm lock}$, exactly as $\Refuse$ filters events to those consistent
with a constraint set.
\end{proposition}

\begin{proof}
The $\Refuse$ operator (Definition in Chapter~\ref{ch:spherepop}) selects events
consistent with $\mathcal{C}$. Here, the events are weight configurations and
$\mathcal{C}_{\rm lock}$ is the constraint that the weight must lie in $\{W_k\}_{k=1}^K$.
Locking maps each detected weight to $\arg\min_k \|W_i - W_k\|$, which is the element
of $\mathcal{C}_{\rm lock}$ nearest to the post-attack weight --- precisely the nearest
consistent configuration. This is a special case of $\Refuse$ under the Euclidean
proximity consistency condition.
\end{proof}

The combined TCU + locking framework therefore implements a two-stage RSVP defense:
the first stage (TCU) engineers the entropy field $\Sfield$ of $\mathcal{M}_W$ to be
low in adversarially vulnerable regions, making large perturbations inadmissible;
the second stage (locking) applies a $\Refuse$-like correction that restores violated
configurations to the nearest admissible centroid.

\begin{theorem}[Admissibility-security duality]
\label{thm:security_duality}
Let $\mathcal{M}_W$ be the weight manifold with RSVP entropy field $\Sfield_W$ encoding
local admissibility radius. Let $HD$ be an adversary's Hamming distance budget and let
$\Sfield_W^{\rm max}(\theta) = \max_{x\in B_{HD}(\theta)}\Sfield_W(x)$ be the maximum
entropy accessible to the adversary within Hamming distance $HD$ of weight $\theta$.
Then the adversary's maximum achievable loss perturbation is bounded by
\begin{equation}
  \max_{HD\text{-attack}} |\Delta\mathcal{L}|
  \leq C \cdot HD \cdot \exp(\Sfield_W^{\rm max}(\theta)),
  \label{eq:security_bound}
\end{equation}
for a constant $C$ depending on the second-order structure of $\mathcal{L}$.
\end{theorem}

\begin{proof}
Each bit-flip of a TCU-encoded weight changes $W_i$ by at most $\Delta_{\rm TCU} =
\exp(-\hat b_i)$ in normalized units. The entropy field $\Sfield_W(\theta)$ at weight
$\theta$ is the log-volume of the admissible perturbation neighborhood, which under
TCU encoding satisfies $\Sfield_W \leq \log(\Delta_{\rm TCU}) = -\hat b_i \log 2$.
Since each of the $HD$ flips contributes at most $C'\exp(\Sfield_W)$ to the loss
perturbation (by the first-order term in~\eqref{eq:onn_sensitivity} and
$|\Delta W^{\rm MSB}_{\rm TCU}| = e^{\Sfield_W}$), summing over $HD$ flips gives
the bound~\eqref{eq:security_bound}.
\end{proof}

Theorem~\ref{thm:security_duality} makes the RSVP security principle precise: the
adversary's damage is bounded by the entropy field of the weight manifold times the
attack budget. Engineering $\Sfield_W$ to be small in sensitive regions is equivalent
to engineering adversarial robustness.

\section{Experimental Grounding}

Lu et al.\ (2024) evaluate the defense framework on VGG-8/CIFAR-10 and ResNet-18/CIFAR-100
with 4-, 6-, and 8-bit quantization under BFA attacks with $HD=100$. Without any defense,
8-bit VGG-8 achieves $13.52\%$ mean post-attack accuracy (from $88\%$ clean accuracy).
The combined TCU + weight locking framework recovers $86.73\%$ mean post-attack accuracy
at $2.36\%$ total memory overhead. This represents a practical realization of the
admissibility-security duality (Theorem~\ref{thm:security_duality}): by engineering
$\Sfield_W$ to be low in the $0.2\%$ most sensitive weights via TCU, and applying
a $\Refuse$-like locking correction post-attack, the framework recovers near-ideal
accuracy at marginal storage cost.

The key quantitative relationships that ground the RSVP interpretation are:

\begin{description}
\item[Sensitivity reduction] TCU reduces MSB-flip perturbation size by $2^{b-1}:1$
  relative to BCD (Proposition~\ref{prop:unary_sensitivity}), which reduces the first-order
  term in sensitivity~\eqref{eq:onn_sensitivity} by the same factor.

\item[Memory compression] TCU achieves $6$--$11\times$ memory reduction over full unary
  (Eq.~\eqref{eq:tcu_bits}), enabling protection of a sufficient fraction $\alpha$
  of vulnerable weights within practical overhead.

\item[Post-attack recovery] Weight locking with $K>1$ centroids outperforms pruning
  ($K=1$, centroid at zero) at equal or lower memory overhead, confirming that
  the $\Refuse$ operation with a richer constraint set $\mathcal{C}_{\rm lock}$
  provides better admissibility recovery.

\item[Noise as partial defense] Hardware Gaussian noise with standard deviation $0.005$
  provides some protection by adding uncertainty to the adversary's gradient estimates,
  but larger noise levels degrade clean accuracy faster than they improve robustness.
  In RSVP terms, noise increases $\Sfield_W$ uniformly (more uncertainty everywhere),
  which helps the defender but also degrades the prior landscape $\Phi$.
\end{description}

\begin{epistbox}
The formal identification of TCU encoding with engineering the RSVP entropy field
$\Sfield_W$, and of weight locking with a $\Refuse$ operation, is an interpretation
within the RSVP framework rather than a claim made by Lu et al.\ (2024). Their paper
establishes the empirical and algorithmic results; the RSVP formalization is the
contribution of this section. Theorem~\ref{thm:security_duality} is a bound derived
within the RSVP framework under assumptions (TCU encoding, first-order loss approximation)
that are validated by the experiments of Lu et al.\ but are stated as approximations
in that work.
\end{epistbox}

\chapter{Topological Defects in RSVP Fields}

\begin{depbox}
This chapter depends on the RSVP field equations~\eqref{eq:rsvp_phi}--\eqref{eq:rsvp_S}
and the admissibility framework (Chapter~\ref{ch:admissibility}).
\end{depbox}

\section{Entropy Vortices and Accessibility Horizons}

A \emph{topological defect} in the RSVP fields is a point (or manifold) where the field
configuration is singular: $\Phi$ diverges, $\Sfield\to 0$, or admissibility fibers
disconnect. These defects are not numerical artefacts; they correspond to physically
meaningful boundaries in the admissibility structure.

The \emph{entropy vortex} is characterized by non-zero circulation
$\Gamma=\oint_C\vfield\cdot d\ell\neq 0$, indicating a persistent loop in the information
flow. Non-zero $\Gamma$ means the admissibility flow does not converge to a fixed point
but circulates---a signature of periodic or quasi-periodic inference dynamics, such as the
oscillatory non-Markovian revivals of Theorem~\ref{thm:non_markov_entropy}.

An \emph{accessibility horizon} is a region where the RSVP velocity field $\vfield$
flows outward in all directions: no trajectory can enter from outside. Inside the horizon,
the entropy field $\Sfield$ may be very low (highly constrained dynamics) or effectively
decoupled from the exterior.

\begin{conjecture}
Long-lived physical and cognitive structures correspond to stabilized entropy defects in
RSVP field geometry: configurations where $\Sfield\approx 0$ in a localized region
(strong constraint, persistent memory) surrounded by $\Sfield>0$ (flexible accessibility).
\end{conjecture}

\begin{exbox}
\textbf{Entropy-vortex simulation.} On a $64\times 64$ lattice with periodic boundary
conditions, initialize $\Phi=\cos(2\pi x/L)\cos(2\pi y/L)$, $\Sfield=0.5+0.3\sin(2\pi x/L)$,
$\vfield=0$. Evolve for $T=200$ steps with $\alpha=0.1$, $\beta=0.05$, $\lambda=0.01$.
Entropy vortices form at the saddle points of $\Phi$ where $\vfield$ circulates around the
$\Sfield$ gradient maximum. Vortex strength (circulation $\Gamma$) grows as $\alpha$ increases.
\end{exbox}

\chapter{RSVP Thermodynamics}

\begin{depbox}
This chapter reinterprets thermodynamic quantities using the RSVP field triple, depending
on Chapters~\ref{ch:admissibility} and the entropy-accessibility identification.
\end{depbox}

\section{Temperature as Accessibility Gradient}

Define the effective temperature in a RSVP system as
\begin{equation}
  T^{-1} = \frac{\partial\Sfield}{\partial E},
\end{equation}
where $E$ is the effective energy of the admissibility configuration. This parallels the
thermodynamic definition $T^{-1}=\partial S/\partial E$ but applies to the accessibility
entropy field rather than thermodynamic entropy.

\section{Entropy Production and the Admissibility Dissipation Functional}

The entropy production rate in the RSVP framework is
\begin{equation}
  \Sigma = \int_\mathcal{M}|\nabla\Sfield|^2\,d\mathrm{vol}_g \geq 0,
\end{equation}
measuring the spatial inhomogeneity of the accessibility field.

\begin{theorem}[Minimum admissibility dissipation]
Steady RSVP states (solutions to $\partial_t\Sfield=0$) minimize $\Sigma$ subject to
fixed boundary conditions on $\Sfield$.
\end{theorem}

\begin{proof}
$\partial_t\Sfield=0$ in equation~\eqref{eq:rsvp_S} gives $\Delta\Sfield\propto\Sfield$ at
leading order (for $\vfield$ given by~\eqref{eq:rsvp_v}). This is the Euler-Lagrange
equation for the functional $\int|\nabla\Sfield|^2$, so steady states minimize the
Dirichlet energy $\int|\nabla\Sfield|^2$, which equals $\Sigma$.
\end{proof}

\chapter{The Geometry of Learning and Semantic Compression}

\begin{depbox}
This chapter connects the RSVP framework to machine learning and representation theory,
depending on the admissibility projection~\eqref{eq:emergent_prob} and the ELBO~\eqref{eq:elbo}.
\end{depbox}

\section{Learning as Accessibility Compression}

Learning is the process of finding a projection $\proj_\theta:\Traj\to\Rep$ that preserves
the admissibility distinctions relevant to a given task while discarding irrelevant ones.
The compression operator is
\begin{equation}
  \mathcal{R}_\theta:\Traj\to\Traj'
\end{equation}
where $\Traj'$ is a reduced trajectory space. Generalization corresponds to the stability
of $\mathcal{R}_\theta$ under perturbations of the training trajectories: if nearby
trajectories in $\Traj$ map to the same region of $\Traj'$, the model generalizes; if
nearby trajectories are separated by $\mathcal{R}_\theta$, the model overfits.

\section{CLIO: Constraint-Leveraged Inference and Optimization}

The CLIO operator acts on the constraint structure rather than on parameters:
\begin{equation}
  \mathrm{CLIO}:(\Traj,\Admiss,\proj)\to(\Traj',\Admiss',\proj').
\end{equation}

A CLIO update penalizes destructive projection and accessibility collapse:
\begin{equation}
  \theta_{t+1} = \theta_t - \eta\nabla_\theta\!\left[
  \mathcal{L}(\proj_\theta(x),d)
  + \lambda\mathcal{D}(\proj_\theta)
  + \mu\mathcal{E}(\Sfield)\right],
\end{equation}
where $\mathcal{D}(\proj_\theta)$ penalizes fiber collapse and $\mathcal{E}(\Sfield)$
penalizes entropy field collapse toward zero.

\begin{implbox}
\textbf{Semantic attractors and transformer embeddings.} Stable regions in representational
geometry correspond to semantic attractors: configurations of $(\Phi,\Sfield)$ where the
RSVP velocity field $\vfield$ converges. In transformer models, token embeddings that are
stable under attention operations correspond to high-$\Phi$, low-$\Sfield$ regions:
high salience, low residual uncertainty. The CLIO penalty $\mathcal{E}(\Sfield)$ prevents
collapse to a single attractor, maintaining diversity of representation.
\end{implbox}

\chapter{Functorial Structure of Spherepop}

\begin{depbox}
This chapter develops the category-theoretic structure of the Spherepop operators, depending
on Chapter~\ref{ch:spherepop} and requiring familiarity with symmetric monoidal categories.
\end{depbox}

\section{Spherepop as a Symmetric Monoidal Category}

Define the \emph{Spherepop category} $\mathbf{Sp}$ with:
objects: event histories $H_t$;
morphisms: admissible transitions $H_t\to H_{t'}$ (causal extensions);
tensor product: $\Bind$ (event composition);
unit: the null event $\bot$.

\begin{proposition}[$\mathbf{Sp}$ is a strict symmetric monoidal category]
The Spherepop category with $\Bind$ as tensor product satisfies the axioms of a strict
symmetric monoidal category: associativity of $\Bind$, existence of unit $\bot$, and
commutativity up to natural isomorphism.
\end{proposition}

\begin{proof}
Associativity: $\Bind(\Bind(e_1,e_2),e_3)=\Bind(e_1,\Bind(e_2,e_3))$ by the associativity
of $\sqcup$ in the event algebra. Unit: $\Bind(e,\bot)=e$ by $e\sqcup\bot=e$. Commutativity
up to isomorphism: $\Bind(e_1,e_2)\cong\Bind(e_2,e_1)$ since $\sqcup$ is commutative
for causally compatible events.
\end{proof}

\section{Collapse as a Functor and Topological Obstructions}

$\Collapse:\mathcal{E}\to\mathcal{H}_{\rm eff}$ is a functor from the event algebra to
the category of effective Hilbert spaces, mapping event composition to tensor product of
compatible subspaces.

\begin{conjecture}
The non-trivial cohomology classes of the Spherepop category $\mathbf{Sp}$ correspond to
irreversible physical memory structures: configurations where the holonomy of the
admissibility connection (Chapter~18) is non-trivial.
\end{conjecture}

\chapter{Computational Complexity of RSVP Inference}

\begin{depbox}
This chapter estimates scaling behavior for RSVP field evolution and KES dynamics,
depending on the field equations~\eqref{eq:rsvp_phi}--\eqref{eq:rsvp_S} and the KES
memory kernel~\eqref{eq:kes_memory}.
\end{depbox}

\section{Scaling of Field Evolution}

\begin{theorem}[Single-step RSVP complexity]
Under bounded entropy gradients $|\nabla\Sfield|\leq C$, a single step of RSVP field
evolution on an $N$-cell lattice (using finite differences) is $O(N)$.
\end{theorem}

\begin{proof}
Each lattice cell requires $O(1)$ operations: computing the Laplacian of $\Phi$ (via
$O(d)$ neighbor values in $d$ dimensions), updating $\vfield$ from~\eqref{eq:rsvp_v},
and updating $\Sfield$ from~\eqref{eq:rsvp_S}. The total is $O(dN)=O(N)$ for fixed $d$.
The entropy gradient bound $|\nabla\Sfield|\leq C$ ensures the nonlinear $\log\Sfield$
term does not cause blow-up.
\end{proof}

\section{Memory Cost and Approximate Inference}

The KES memory kernel $K_m$ of depth $m$ requires storing the last $m$ event vectors:
$O(m|\mathcal{E}|)$ memory. For deep memory ($m\gg 1$), variational RSVP approximates
the full kernel by a parametric family:
\begin{equation}
  K_m^\phi(e_{t-m},\ldots,e_{t-1}) \approx f_\phi(e_{t-m},\ldots,e_{t-1}),
\end{equation}
where $f_\phi$ is a neural network trained to minimize the KL divergence from the true
kernel. This gives $O(|\phi|)$ memory independent of $m$, at the cost of approximation error.

\begin{implbox}
\textbf{GPU parallelism for RSVP evolution.} Each RSVP lattice cell is an independent
CUDA thread. Fields $(\Phi_i,\vfield_i,\Sfield_i)$ are stored as structure-of-arrays in
global memory; the update kernel~\eqref{eq:rsvp_phi}--\eqref{eq:rsvp_S} is launched with
$N$ threads per step. Asynchronous entropy propagation uses shared memory for the
$\operatorname{div}(\Sfield\vfield)$ term within each warp.

\medskip
\textbf{Complexity summary.}
\begin{center}
\begin{tabular}{@{}lll@{}}
\toprule
Component & Per-step cost & Memory \\
\midrule
RSVP field evolution & $O(N)$ & $O(N)$ \\
KES event selection & $O(|\mathcal{E}|)$ & $O(|\mathcal{X}_t|)$ \\
KES memory ($m$ steps) & $O(m|\mathcal{E}|)$ & $O(m|\mathcal{E}|)$ \\
Variational KES & $O(|\phi|)$ & $O(|\phi|)$ \\
TARTAN ($K$ cells) & $O(KN)$ & $O(KN)$ \\
\bottomrule
\end{tabular}
\end{center}
\end{implbox}

\chapter{RSVP Cosmological Optics}

\begin{depbox}
This chapter applies the RSVP propagator~\eqref{eq:rsvp_propagator} to cosmological photon
propagation. It is conjectural throughout; epistemic status is stated explicitly.
\end{depbox}

\begin{epistbox}
This chapter presents conjectural cosmological applications of RSVP rather than
experimentally established physics. All results should be understood as predictions of
the framework contingent on its physical identification with the photon propagation equation.
\end{epistbox}

\section{Entropic Redshift}

The RSVP propagator predicts a cosmological photon frequency shift
\begin{equation}
  \frac{\Delta\omega}{\omega_0} = \frac{1}{2}(\Phi(x_{\rm obs})-\Phi(x_{\rm emit}))
  - \frac{\alpha}{2}\int_{\rm path}\Sfield\,ds.
\end{equation}

If $\Phi$ decreases along the line of sight (earlier epochs have lower likelihood potential),
this produces a redshift without metric expansion. The entropic attenuation term modifies
the luminosity-distance relation relative to $\Lambda$CDM.

\section{Structure Formation as Accessibility Condensation}

Large-scale structure forms where the RSVP scalar field $\Phi$ develops local maxima
(high admissibility basins): matter and energy are drawn toward high-$\Phi$ regions by
the RSVP velocity field $\vfield=(\alpha/\Sfield)\nabla\Phi$. This replaces gravitational
collapse as the primary mechanism of structure formation, with the gravitational potential
identified as a component of $\Phi$.

\begin{conjecture}
Observed cosmological acceleration may arise from large-scale entropy-gradient effects
($\nabla\Sfield\neq 0$ on Hubble scales) rather than from a cosmological constant.
\end{conjecture}

% =============================================================
\part{Toward a Unified Accessibility Calculus}
% =============================================================

\chapter{Projection, Agency, and Observer Geometry}

\begin{depbox}
This chapter depends on the admissibility framework (Chapter~\ref{ch:admissibility}) and
the Spherepop axioms (Chapter~\ref{ch:spherepop}).
\end{depbox}

\section{Observers as Recursive Projection Systems}

An observer is a system that recursively projects $\Traj$ onto a sequence of increasingly
compressed representations $\Rep_1\supset\Rep_2\supset\cdots$, retaining only admissible
distinctions at each level. Define the \emph{observer projection map}
$\proj_i:\Traj\to\Rep_i$, where $\Rep_i$ is the $i$-th observer's representational manifold.

\begin{definition}[Consensus manifold]
The \emph{consensus manifold} for a set of observers $\{\proj_i\}$ is
\begin{equation}
  \Rep_{\rm consensus} = \bigcap_i\proj_i(\Traj),
\end{equation}
the intersection of representational images restricted to mutually stable admissibility regions.
\end{definition}

Shared classical reality corresponds to a large, stable consensus manifold. Quantum
non-classicality corresponds to configurations where different observer projections disagree:
$\proj_i(x)\neq\proj_j(x)$ for trajectories $x$ in the non-classical regime.

\begin{conjecture}
Persistent observer identity corresponds to a stable fixed point of recursive admissibility
compression: a configuration $\proj^*$ such that $\proj^*\circ\proj^*\approx\proj^*$ in
the sense that repeated self-projection does not substantially reduce the admissibility structure.
\end{conjecture}

\chapter{The Unified Accessibility Calculus and Open Problems}

\begin{depbox}
This final chapter synthesizes all preceding material and defines the long-term research
program. All results marked as conjectures or open problems are formally stated but unproved.
\end{depbox}

\section{The Universal Accessibility Functional}

Define the \emph{total admissibility} of a system at time $t$:
\begin{equation}
  \mathfrak{A}(t) = \int_\Traj\Admiss(x,t)\,dx.
\end{equation}

The unified evolution equation for the admissibility functional is proposed as
\begin{equation}
  \partial_t\Admiss = \mathcal{G}(\Phi,\vfield,\Sfield,H_t),
\end{equation}
where $\mathcal{G}$ is a generalized admissibility operator encoding the coupled RSVP
dynamics, the KES event selection, and the history-dependent memory kernel. The RSVP
field equations~\eqref{eq:rsvp_phi}--\eqref{eq:rsvp_S} are the leading-order approximation
to $\mathcal{G}$ in the limit of small history dependence ($m=0$) and compact manifold.

\section{Mathematical Open Problems}

\begin{description}
\item[Global existence.] Theorem~\ref{thm:rsvp_wellposed} gives local-in-time well-posedness.
  Whether solutions exist globally or develop finite-time singularities (entropy field reaching
  zero) remains open. The lower bound $\Sfield\geq s_0 e^{-Ct}$ suggests sufficiently strong
  $\beta$ is needed.

\item[Categorical coherence of $\mathbf{Sp}$.] The Spherepop category $\mathbf{Sp}$ should
  admit a dagger structure encoding time-reversal asymmetry. Whether $\mathbf{Sp}$ embeds
  into the Abramsky-Coecke categorical framework for quantum mechanics is open.

\item[Fiber bundle curvature and Berry phase.] The conjecture that RSVP admissibility
  connection holonomy equals the quantum mechanical Berry phase requires a formal identification
  of the RSVP horizontal distribution with the quantum geometric tensor.

\item[Quantization of RSVP.] The classical field theory $\mathcal{L}_{\rm RSVP}$ should be
  quantized. The positivity of $\Sfield$ in the classical theory must be preserved as a
  constraint on the quantum Hilbert space.

\item[Complexity of exact KES.] The exact KES map requires solving the optimization
  $\max_e\mathcal{F}(e)$ over $\mathcal{E}$. The computational complexity of this
  optimization as a function of $|\mathcal{E}|$ and $|\mathcal{X}_t|$ is not known in general.
\end{description}

\section{The Four Core Ideas}

Every open problem, extension, and application in this text is anchored to the same four
conceptual pillars, which the reader should hold in view throughout:

\begin{enumerate}
\item \textbf{Admissibility}: not all configurations are equally accessible; differential
  accessibility is the true primitive.

\item \textbf{Projection}: the map $\proj:\Traj\to\Rep$ and its fiber structure encode
  the information cost of conventional modeling.

\item \textbf{Entropy as accessibility}: $\Sfield$ measures future trajectory volume,
  not merely Shannon uncertainty.

\item \textbf{Irreversible event structure}: events are primitive; states are equivalence
  classes of event histories; irreversibility is foundational, not emergent.
\end{enumerate}

These four ideas are the invariants of the framework. Any section that cannot be connected
to at least two of them is a candidate for removal.

\backmatter

\bibliographystyle{plainnat}
\begin{thebibliography}{99}

\bibitem[Amari(2016)]{amari2016}
Amari, S. (2016). \textit{Information Geometry and Its Applications}. Springer.

\bibitem[Lu et al.(2024)]{lu2024}
Lu, H., Yin, Z., Bhoumik, P., Banerjee, S., Chakrabarty, K., and Gu, J. (2024).
The unlikely hero: Nonideality in analog photonic neural networks as built-in defender
against adversarial attacks. \textit{arXiv:2410.01289}. To appear in \textit{Proc.\ ASP-DAC}, 2025.

\bibitem[Born(1926)]{born1926}
Born, M. (1926). Zur Quantenmechanik der Stossvorg\"{a}nge.
\textit{Zeitschrift f\"{u}r Physik}, 37, 863--867.

\bibitem[Caticha(2019)]{caticha2019}
Caticha, A. (2019). The entropic dynamics approach to quantum mechanics.
\textit{Entropy}, 21(10), 943.

\bibitem[Fuchs, Mermin and Schack(2014)]{fuchs2014}
Fuchs, C.~A., Mermin, N.~D., and Schack, R. (2014).
An introduction to QBism. \textit{American Journal of Physics}, 82(8), 749--754.

\bibitem[Friston(2010)]{friston2010}
Friston, K. (2010). The free-energy principle: a unified brain theory?
\textit{Nature Reviews Neuroscience}, 11(2), 127--138.

\bibitem[Gelman et al.(2013)]{gelman2013}
Gelman, A., Carlin, J.~B., Stern, H.~S., Dunson, D.~B., Vehtari, A., and Rubin, D.~B. (2013).
\textit{Bayesian Data Analysis}. 3rd ed. CRC Press.

\bibitem[Gleason(1957)]{gleason1957}
Gleason, A.~M. (1957). Measures on the closed subspaces of a Hilbert space.
\textit{Journal of Mathematics and Mechanics}, 6(6), 885--893.

\bibitem[Gorini, Kossakowski and Sudarshan(1976)]{gks1976}
Gorini, V., Kossakowski, A., and Sudarshan, E.~C.~G. (1976).
Completely positive dynamical semigroups of $N$-level systems.
\textit{Journal of Mathematical Physics}, 17(5), 821--825.

\bibitem[Griffiths(1984)]{griffiths1984}
Griffiths, R.~B. (1984). Consistent histories and the interpretation of quantum mechanics.
\textit{Journal of Statistical Physics}, 36(1), 219--272.

\bibitem[Hong, Ou and Mandel(1987)]{hom1987}
Hong, C.~K., Ou, Z.~Y., and Mandel, L. (1987).
Measurement of subpicosecond time intervals between two photons by interference.
\textit{Physical Review Letters}, 59(18), 2044.

\bibitem[Jaynes(1957)]{jaynes1957}
Jaynes, E.~T. (1957). Information theory and statistical mechanics.
\textit{Physical Review}, 106(4), 620.

\bibitem[Kasap(2013)]{kasap2013}
Kasap, S.~O. (2013).
\textit{Optoelectronics and Photonics: Principles and Practices}. 2nd ed. Pearson.

\bibitem[Lindblad(1976)]{lindblad1976}
Lindblad, G. (1976). On the generators of quantum dynamical semigroups.
\textit{Communications in Mathematical Physics}, 48(2), 119--130.

\bibitem[Mandel and Wolf(1995)]{mandel1995}
Mandel, L. and Wolf, E. (1995).
\textit{Optical Coherence and Quantum Optics}. Cambridge University Press.

\bibitem[Rivas and Huelga(2012)]{rivas2012}
Rivas, A. and Huelga, S.~F. (2012).
\textit{Open Quantum Systems: An Introduction}. Springer.

\bibitem[Rovelli(1996)]{rovelli1996}
Rovelli, C. (1996). Relational quantum mechanics.
\textit{International Journal of Theoretical Physics}, 35(8), 1637--1678.

\bibitem[Schlosshauer(2007)]{schlosshauer2007}
Schlosshauer, M. (2007).
\textit{Decoherence and the Quantum-to-Classical Transition}. Springer.

\bibitem[Soshnikov(1999)]{soshnikov1999}
Soshnikov, A. (1999). Universality at the edge of the spectrum in Wigner random matrices.
\textit{Communications in Mathematical Physics}, 207(3), 697--733.

\bibitem[Tracy and Widom(1994)]{tracy_widom1994}
Tracy, C.~A. and Widom, H. (1994). Level-spacing distributions and the Airy kernel.
\textit{Communications in Mathematical Physics}, 159(1), 151--174.

\bibitem[Wakefield(2013)]{wakefield2013}
Wakefield, J. (2013).
\textit{Bayesian and Frequentist Regression Methods}. Springer.

\bibitem[Zurek(2003)]{zurek2003}
Zurek, W.~H. (2003). Decoherence, einselection, and the quantum origins of the classical.
\textit{Reviews of Modern Physics}, 75(3), 715.

\end{thebibliography}

\printindex
\end{document}
